%\chapter{Reference guide}

\section{Reference guide}

\subsection{Units and variables}

The following set of units, based on SI with some exceptions, is adopted
in the Astra code. 
\medskip\\
\noindent
{\bf Table 4.1. System of units}
\nopagebreak
\\[10pt]
\nopagebreak
\noindent
\begin{tabular}{|l|l|}
\hline
\bf Unit		& \bf Quantity\\
\hline
m			& length\\
m$^2$			& area\\
m$^3$			& volume\\
s			& time\\
m/s			& velocity\\
m$^2$/s			& diffusion\\
$10^{19}$m$^{-3}$	& particle density\\
$10^{19}$m$^{-3}$s$^{-1}$ & particle source\\
$10^{19}$s$^{-1}$	& particle flux\\
$10^{19}$m$^{-2}$s$^{-1}$ & particle flux density\\
keV			& temperature\\
s$^{-1}$		& collision frequency\\
\hline
\end{tabular}
\hspace*{30pt}\begin{tabular}{|l|l|}
\hline
\bf Units		& \bf Quantity\\
\hline
MW			& power, energy flux\\
MW$\cdot$m$^{-3}$	& power density\\
MW$\cdot$m$^{-2}$	& density of energy flux\\
MJ			& energy content\\
T			& magnetic induction\\
V$\cdot$s		& magnetic flux\\
V			& voltage\\
MA			& current\\
MA$\cdot$m$^{-2}$	& current density\\
$\mu$Hn			& inductance\\
($\mu$Ohm$\cdot$m)$^{-1}\!=$MS/m	& conductivity\\
GHz			& RF frequency\\
\hline
\end{tabular}
\\[10pt]
The main set of the simple variables in the Astra code includes
\medskip\\
\noindent
{\bf Table 4.2. List of main device and plasma parameters}
\nopagebreak
\\[10pt]
\nopagebreak
\noindent
\begin{tabular}{|c|l|l|l|}
\hline
\bf Quantity	&\bf Variable name &\bf Units &\bf Description	\\
\hline
$R_0$		&{\tt RTOR}	& m   & major radius		 \\
$a_M$		&{\tt AB}	& m   & maximum value of $a_B$ allowed \\
$a_B$		&{\tt ABC}	& m   
			& minor radius of LCMS in the mid-plane \\
$\Delta_B$	&{\tt SHIFT}	& m   
			& Shafranov shift of the plasma boundary\\
$\lambda_B$	&{\tt ELONG}	& -- 
			& elongation of the plasma boundary \\
$\delta$	&{\tt TRIAN}	& -- 
			& triangularity of the plasma boundary \\
$\Delta_Z$	&{\tt UPDWN}	& m   
			& vertical shift of the plasma boundary	 \\
$\rho_B$	&{\tt ROC}	& m   & effective minor radius	 \\
$A_i/m_p$	&{\tt AMJ}	&--
			& main ion mass number ($m_p$ is the proton mass)\\
$Z_i$		&{\tt ZMJ}	& -- 
			& electric charge of the main ion species \\
$B_0$ 		&{\tt BTOR}	& T   & vacuum magnetic field at $R_0$ \\
$I_{\rm pl}$	&{\tt IPL}	& MA  & total plasma current 	\\
%$L_{\rm ext}$ 	&{\tt LEXT}	& $\mu$Hn&external inductance	 \\
%$U_{\rm ext}$ 	&{\tt UEXT}	& V   & external voltage	 \\
%$V$ 		&{\tt VOLUME}	& m$^3$   & plasma volume	 \\
%$t$ 		&{\tt TIME}	& s   & time			 \\
%$\tau$	 	&{\tt TAU}	& s   & time step		 \\
%$\pi$		&{\tt GP}	& -- & 3.1415926		 \\
\hline
\end{tabular}
\\[10pt]
The first two quantities in this table cannot depend on time, all others 
can be time dependent.

Table 4.3 shows the list of main radially and time dependent
quantities (arrays) which are continuously stored in memory at two
time slices.  
Except for $n_i$ and $V',$ this list coincides with the list of unknown
quantities in the system of transport equations. 
\medskip\\
\noindent
{\bf Table 4.3. List of radial functions (arrays) stored at two time
slices} 
\nopagebreak
\\[10pt]
\nopagebreak
\noindent
\begin{tabular}{|c|c|c|l|}
\hline
\bf Quantity	&\bf Code variable at the
		&\bf Code variable at the 	&\bf Description\\ 
	&\bf current time $t\hspace*{2.6em}$ 
		&\bf previous time $t-\tau\hspace*{0.9em}$&	\\ \hline
$n_e$ 		&{\tt NE}	&{\tt NEO}	& electron density\\
$n_i$		&{\tt NI}	&{\tt NIO}	& ion density\\
$T_e$		&{\tt TE}	&{\tt TEO}	& electron temperature\\
$T_i$ 		&{\tt TI}	&{\tt TIO}	& ion temperature\\
$\psi$ 		&{\tt FP}	&{\tt FPO}	& poloidal flux\\
$f_1$ 		&{\tt F1}	&{\tt F1O} & 1st function in \req{aux}\\
$f_2$ 		&{\tt F2}	&{\tt F2O} & 2nd function in \req{aux}\\
$f_3$ 		&{\tt F3}	&{\tt F3O} & 3rd function in \req{aux}\\
$V'$ 		&{\tt VR}	&{\tt VRO} & volume derivative\\ \hline
\end{tabular}
\\[10pt]
All other quantities are available for the current time, $t,$ only. 
The list of main arrays is given in Table 4.4. 
\medskip\\
\noindent
{\bf Table 4.4. Some other radial functions}
\nopagebreak
\\[10pt]
\nopagebreak
\noindent
\begin{tabular}{|c|l|l|l|}
\hline
\bf Quantity	&\bf Array name	&\bf Units	&\bf Description\\
\hline
$\rho$		&{\tt RHO}	& m	& main magnetic surface label\\
$a$		&{\tt AMETR}	& m	& minor radius as in \req{3mom}\\
$\Delta$	&{\tt SHIF}	& m	& Shafranov shift \req{3mom}\\
$\lambda$	&{\tt ELON}	& --	& elongation \req{3mom}\\
$\delta$	&{\tt TRIA}	& --	& triangularity \req{3mom}\\
%$n_i$ 		&{\tt NI}	& $10^{19}$m$^{-3}$ & ion density\\
$j_\pll$ 	&{\tt CU}	& A/m$^2$ & current density\\
$j_{BS}$ 	&{\tt CUBS}	& A/m$^2$ & bootstrap current density\\
$j_{CD}$ 	&{\tt CD}	& A/m$^2$ & driven current density\\
$\sigma_\pll$ 	&{\tt CC} & ($\mu$Ohm$\cdot$m)$^{-1}$	& conductivity\\
$\mu$	 	&{\tt MU}	& --	& rotational transform\\
$\mu_V$		&{\tt MV}	& --	& vacuum rotational transform\\
$\psi_V$	&{\tt FV}	& V$\cdot$s
			& poloidal flux for vacuum magnetic field\\
$E_\pll/B_p$	&{\tt VP}	& m/s	& pinch velocity, 
					$E_\pll=U_\pll/(2\pi R_0)$\\
$U_\pll$	&{\tt ULON}	& V	& longitudinal loop voltage\\
$U_{\rm pl}$	&{\tt UPL}	& V	& toroidal loop voltage\\
$Z_{eff}$ 	&{\tt ZEF}	& --	& effective charge\\
$\Gamma_e$ 	&{\tt QN}	& $10^{19}$s$^{-1}$ & particle flux\\
$q_e$	 	&{\tt QE}	& MW	& electron heat flux\\
$q_i$	 	&{\tt QI}	& MW	& ion heat flux\\
$J$ 		&{\tt IPOL}	& --	& normalized poloidal current\\
$V'G_1$		&{\tt G11}	& m$^2$	& defined in \req{metr}\\
$R_0G_2/J$	&{\tt G22}	& m	& defined in \req{metr}\\
$G_3$		&{\tt G33}	& --	& defined in \req{metr}\\
\hline
\end{tabular}
\pagebreak
\\
The complete list of radially dependent variables includes more than 200
quantities. 
Most of them are discussed in Section 4.6.2.

\subsection{Transport equations in Astra notations}

For convenience in working with the Astra code, we re-write the set of
transport equations (\ref{main})-(\ref{fluxes}) in the newly introduced
system of units using Astra variable notation. First, we define
the operators 
\begin{equation}
\begin{array}{c}
\hat D_T[f]\df\dstyle
\frac{3}{2}\frac{1}{\lb(V'\rb)^{5/3}}
\lb(\prti{}{t}-\frac{\dot B_0}{2B_0}\prti{}{\rho}\rho\rb)
\lb[\lb(V'\rb)^{5/3}f\rb]
,\bigskip\\ \dstyle
\hat D_n[f]\df
\frac{1}{V'}\lb(\prti{}{t}-\frac{\dot B_0}{2B_0}\prti{}{\rho}\rho\rb)
\lb[V'f\rb]
,\qquad\dstyle
\hat D_\psi[f]\df
\prti{f}{t}-\frac{\rho\dot B_0}{2B_0}\prti{f}{\rho}
.\end{array}
\end{equation}
Now \req{main} takes the form
\begin{equation}\label{Amain}
\hspace*{-5mm}
\begin{array}{rl}\dstyle
\hat D_n\left[{\tt NE}\right]+\prti{}{V}\lb({\tt QN}\rb)
&\!\!\!\!
={\tt SN}+{\tt SNN*NE}
,\bigskip\\ \dstyle
\hat D_T\left[{\tt NE*TE}\right]
+\prti{}{V}
\lb({\tt 625*QE}+\frac{\tt 5}{\tt 2}{\tt *}\gamma_e{\tt *TE*QN}\rb)
&\!\!\!\!
={\tt 625*(PE+PET*TE)}
,\bigskip\\ \dstyle
\hat D_T\left[{\tt NI*TI}\right]
+\prti{}{V}
\lb({\tt 625*QI}+\frac{\tt 5}{\tt 2}{\tt*}\gamma_i\;
\frac{\tt NI*TI}{\tt NE}\;{\tt QN}\rb)
&\!\!\!\!
={\tt 625*(PI+PIT*TI)}
,\bigskip\\ \dstyle
{2\pi\rho}{\tt *CC*}\hat D_\psi\lb[{\tt FP}\rb]
+V'{\tt *}\left({\tt CUBS}+{\tt CD}\right)
&\!\!\!\!\dstyle
={\tt 5*IPOL}^2\prti{}{\rho}
\left[{\tt G22}\prti{}{\rho}\lb(\tt FP\rb)\right]
.\end{array}
\end{equation}
Two additional factors $\gamma_e,\gamma_i$ are introduced in \req{Amain}
to allow modification of the contribution of the convective component to
the heat flux. 
The corresponding code variable names are
$\tw{GN2E}=\frac{5}{2}\gamma_e$ and $\tw{GN2I}=\frac{5}{2}\gamma_i.$
When not determined explicitly these factors take the default values 
$\tw{GN2E}=\tw{GN2I}=0.$ 
The full energy fluxes through a whole magnetic surface appear in the
form 
\begin{equation}
Q_e=q_e+\frac{5}{2}\gamma_e T_e \Gamma_e
={\tt QE}+1.6\x10^{-3}\,{\tt GN2E*TE*QN}
\end{equation}
and
\begin{equation}
Q_i=q_i+\frac{5}{2}\gamma_i T_i \Gamma_i
={\tt QI}+1.6\x10^{-3}\,\frac{\tt NI}{\tt NE}\;{\tt GN2I*TI*QN}
\end{equation}
where the factor $1.6\x10^{-3}=1/625$ emerges because the mixed system of units (keV and
MJ) is adopted. 

The set of equations (\ref{aux}) acquires the form
\begin{equation}\label{Aux}
\frac{1}{V'}\prti{}{t}\left(V'{\tt *Fj}\right)
+\prti{}{V}\lb({\tt QFj}\rb)
={\tt SFj+SFFj*Fj}
,\end{equation}
with
\begin{equation}\label{QFj}
{\tt QFj}
={\tt G11*}
\left({\tt VFj*Fj}-{\tt DFj}\;\prti{}{\rho}\lb({\tt Fj}\rb)\right)
.\end{equation}
Here the letter {\tt j} should be replaced with one of the numbers: 1,
2, 3. 

The transport matrix \req{fluxes}, in the Astra notation, takes the form
\begin{equation}\label{Afluxes}
\left(
\begin{array}{c}\dstyle
\frac{\tt QN-SLAT*GNX}{\tt G11*NE}
\bigskip\\ \dstyle
\frac{\tt 625*QE}{\tt G11*NE*TE}
\bigskip\\ \dstyle
\frac{\tt 625*QI}{\tt G11*NI*TI}
\bigskip\\ \dstyle
\frac{0.4\pi}{B_p}{\tt *CUBS}
\end{array}
\right)
=
-\left(
\begin{array}{cccc}
\tt DN &\tt HN &\tt XN &\tt CN
\bigskip\bigskip\\ \dstyle
\tt DE &\tt HE &\tt XE &\tt CE
\bigskip\bigskip\\ \dstyle
\tt DI &\tt HI &\tt XI &\tt CI
\bigskip\bigskip\\ \dstyle
\tt DC &\tt HC &\tt XC & 0
\end{array}
\right)
\cdot
\left(
\begin{array}{c}\dstyle
\frac{1}{\tt NE}\prti{\tt NE}{\rho}
\bigskip\\ \dstyle
\frac{1}{\tt TE}\prti{\tt TE}{\rho}
\bigskip\\ \dstyle
\frac{1}{\tt TI}\prti{\tt TI}{\rho}
\bigskip\\ \dstyle
-1
\end{array}
\right)
.\end{equation}
The reason for the additional term \tws{-SLAT*GNX} in the first
row on the left hand side of \req{Afluxes} will be discussed in
Section~4.9.1. 
Note also that the last 
element of the column vector on the \rhs of \req{fluxes} does not
correspond to that in \req{Afluxes}:  
namely, the quantity ``$E_\pll/B_p$'' in \req{fluxes} is replaced with
``1'' in \req{Afluxes}. 
This allows the inclusion of particle or energy fluxes which are not
necessarily related to the toroidal electric field. 
\medskip\\
\noindent
{\bf Table 4.5. Notations for transport coefficients in the Astra code}
\nopagebreak
\\[10pt]
\nopagebreak
\noindent
\begin{tabular}{| l l l | l | l l l|}				 \hline
\multicolumn{3}{|c|}{\bf Code variable}	&\bf Units 	&
\multicolumn{3}{|l|}{Notations in Eqs.~(\ref{fluxes}),(\ref{aux})}\\
\hline
{\tt DN,}&{\tt HN,}&{\tt XN}    &m/s$^2$ 
		&\qquad $D_n$,       &$D_e$,       &$D_i$	\\
{\tt DE,}&{\tt HE,}&{\tt XE}    &m/s$^2$ 
		&\qquad $\chi_n^e,$&$\chi_e,$  &$\chi_i^e$\\
{\tt DI,}&{\tt HI,}&{\tt XI}    &m/s$^2$ 
		&\qquad $\chi_n^i,$&$\chi_e^i,$&$\chi_i$	\\
{\tt CN,}&{\tt CE,}&{\tt CI}    &m/s     
		&\qquad $C_n,$       &$C_e,$       &$C_i$ 	\\
{\tt DC,}&{\tt HC,}&{\tt XC}    &--     
		&\qquad $D_E,$       &$\chi_E^e,$&$\chi_E^i$\\
{\tt DF1,}&{\tt DF2,}&{\tt DF3} &m/s$^2$ 
		&\qquad $D_1,$       &$D_2,$       &$D_3$	\\
{\tt VF1,}&{\tt VF2,}&{\tt VF3} &m/s     
		&\qquad $v_1,$       &$v_2,$       &$v_3$	\\
{\tt CC} & & 		&MS/m	&\qquad $\sigma_\pll$ & &	\\[5pt]
\hline
\end{tabular}\vskip 3ex

The sources on the right hand side in Eqs.~(\ref{Amain}) and (\ref{Aux})
are split into two parts. 
The part which is linear with respect to the unknown variable in each
equation is written separately, in order to allow the choice of either
an implicit or an explicit numerical scheme.
For instance, the same particle source $S_e$ can be defined as 
${\tt SN}=s_{ion}Nn_e$ or as ${\tt SNN}=s_{ion}N.$
In the first case, it would be approximated explicitly using the
electron density $n_e={\tt NEO}$ from the previous time slice,
$S_e=s_{ion}N*{\tt NEO};$
while, in the second case, $S_e=s_{ion}N*{\tt NE}.$
We summarize usage of the sources and sinks in Table 4.6.

\vskip 1ex\noindent
{\bf Table 4.6. Notations for the right hand sides of transport equations}
\\[10pt]\noindent
\begin{tabular}{| lll | l | l |}				\hline
\multicolumn{3}{|c|}{\bf Code variable}
			&\bf Units	&\bf Definition\\	\hline
{\tt SN}  & 	&    &$10^{19}$/(m$^3\,$s)& $S_e={\tt SN+SNN*NE}$	\\
{\tt SNN} & 	&	&1/s		&		 	\\
{\tt PE,} &{\tt PI,} &	&MW/m$^3$	& $P_e={\tt PE+PET*TE}$	\\
{\tt PET,}&{\tt PIT,}&	&MW/(m$^3\,$keV)	& $P_i={\tt PI+PIT*TI}$	\\
{\tt SF1,}&{\tt SF2,}&{\tt SF3} &[{\tt Fj}]/s & $S_j={\tt SFj+SFFj*Fj}$\\
{\tt SFF1,}&{\tt SFF2,}&{\tt SFF3}&1/s	& 			\\
{\tt CUBS,}&{\tt CD}&	&MA/m$^2$	& $j_{BS}={\tt CUBS},$
					  $j_{CD}={\tt CD}$	\\ \hline
\end{tabular}

\vskip 3ex
The transport set of equations will be completely defined if 
\begin{description}
\vspace*{-1ex}
\item (1) device and plasma parameters (Table 4.2) are set,
\vspace*{-1ex}
\item (2) all quantities in Tables 4.5 and 4.6 are determined,
\vspace*{-1ex}
\item (3) a rule for calculation of metric coefficients is defined and 
\vspace*{-1ex}
\item (4) initial and boundary conditions are prescribed. 
\end{description}
Calculation of the metric coefficients is provided by the equilibrium
solver, and 
a link between transport and equilibrium parts of the code is discussed
in Section 4.5.
The technical details for setting initial and boundary conditions
in the code are described in Section 5, while in the following 
two sections the general logic of the procedure is discussed.

\subsection{Initial conditions}

In this Section and below we use the following flowchart notation
$$
\begin{picture}(110,50)
\put(0,30){\begin{picture}(110,50)
	\thicklines
	\multiput(15,15)(0,-30){2}{\line(1,0){30}}
	\multiput(0,0)(45,-15){2}{\line(1,1){15}}
	\multiput(0,0)(45,15){2}{\line(1,-1){15}}
	\put(20,-5){{\tt QTY}}
	\thinlines
	\put(67,0){\vector(1,0){40}}\put(80,5){yes}
	\put(30,-20){\vector(0,-1){25}}\put(15,-32){no}
	\end{picture}}
\end{picture}\QQUAD
$$
which means a test of whether the quantity {\tt QTY} is defined in a
model or not. 
If it appears on the left hand side of any assignment instruction in 
a model (see Section 5.2) at least once as
\\ \indent
{\tt QTY}$=${\tt expression};\\
then we say that it is defined, the answer to the question is positive 
and the path labeled "yes" is taken.


\subsubsection{Initial conditions for $\tt NE, TE, TI, Fj$}

Setting of the initial conditions for the electron density is quite
straightforward:\\ 
\hspace*{125pt}\begin{tabular}{cl}
\begin{picture}(110,25)
\thicklines
\multiput(15,15)(0,-30){2}{\line(1,0){30}}
\multiput(0,0)(45,-15){2}{\line(1,1){15}}
\multiput(0,0)(45,15){2}{\line(1,-1){15}}
\put(25,-5){{\tt NE}}
\thinlines
\put(67,0){\vector(1,0){40}}\put(80,5){yes}
\put(30,-20){\vector(0,-1){30}}\put(15,-35){no}
\end{picture}
&
$n_{e}(\rho,t)={\tt NE}$\\[50pt]
%\begin{picture}(110,65)
%\thicklines
%\multiput(15,15)(0,-30){2}{\line(1,0){30}}
%\multiput(0,0)(45,-15){2}{\line(1,1){15}}
%\multiput(0,0)(45,15){2}{\line(1,-1){15}}
%\put(20,-5){{\tt NEX}}
%\thinlines
%\put(70,0){\vector(1,0){40}}\put(80,5){yes}
%\put(30,-20){\vector(0,-1){30}}\put(15,-35){no}
%\end{picture}
%&
\hspace*{-50pt}$n_{e}(\rho,t)={\tt NEX}$\\
%\multicolumn{2}{l}
%{\hspace*{-50pt}The density distribution is determined}\\ 
%\multicolumn{2}{l}
%{\hspace*{-50pt}by the default value $10^{18}\rm\;m^{-3}$}\\
\end{tabular}\\[15pt]
Firstly, the Astra compiler checks whether the diffusion equation
for $n_e$ is to be solved or not.
Suppose first that the density evolution will be prescribed, i.e. the
particle diffusion equation is to be suppressed. 
The Astra compiler continues by checking whether \twd{NE} is assigned
in the model. 
If yes, then the density will evolve according to this assignment and no
further checks are made.
If \twd{NE} is not defined in the model, then the next check (not shown
in the diagram) is made as to whether \twd{NEX} is set in the data file
(Section~5.4). 
If this is the case, then \twd{NE} is set to \twd{NEX}.
If it happens that \twd{NEX} does not appear in the data file either,
then the default setting $n_e(\rho,t)=10^{18}{\rm\;m}^{-3}$ is applied. 
The default setting is constant in time and space, while all other means
of setting \twd{NE} can have arbitrary dependences.

A similar scheme is applied when the density $n_e(\rho,t)$ is to be
defined by the transport equation.
The difference is that, in this case, the assignment
$n_e(\rho,t_0)=\dots$ is made only once, and that this is understood as
the initial condition. 
Any time $t_0$ (code variable \twd{TSTART}) can be selected as a starting
time for the simulation. 
The following evolution of $n_e$ at $t>t_0$ will be defined by the
transport equation. 
Finally, it has to be mentioned that there is an option (see
Section~5.2.5) to assign \twd{NE} only once at $t=t_0=\tt TSTART$ even
if \twd{NE} or \twd{NEX} are defined as time dependent. 

Initial conditions for $T_e,$ $T_i$ and $f_j$ are defined in similar
ways:
\pagebreak 

%\indent
{{\bf Table 4.7. Summary of the  setting of initial conditions}
\bigskip\\ \noindent
\begin{tabular}{| c | c | c | c | l |}\hline
{\bf Variable}	&{\bf Name in the code}
	&\bf 1st check  &\bf 2nd check	&\bf Default value \\ \hline
$n_e$	& {\tt NE} &${\tt NE}=\:?$ & ${\tt NEX}=\:?$
				   & \quad $10^{18}{\rm m}^{-3}$ \\
$T_e$	& {\tt TE} &${\tt TE}=\:?$ & ${\tt TEX}=\:?$ & \quad 10 eV \\
$T_i$	& {\tt TI} &${\tt TI}=\:?$ & ${\tt TIX}=\:?$ & \quad 10 eV \\
$f_1$	& {\tt F1} &${\tt F1}=\:?$ & ${\tt F1X}=\:?$ & \quad 0	\\
$f_2$	& {\tt F2} &${\tt F2}=\:?$ & ${\tt F2X}=\:?$ & \quad 0	\\
$f_3$	& {\tt F3} &${\tt F3}=\:?$ & ${\tt F3X}=\:?$ & \quad 0	\\
\hline
\end{tabular}}

\vskip 2ex\noindent
For every quantity in this table, the Astra compiler finds the
appropriate mode of calculating the corresponding variable:
via equation, formula or data file.

\subsubsection{Initial conditions for the poloidal flux $\psi=\tt FP$}

The similar procedure is applied for the poloidal flux.
The scheme below can define the initial conditions only (then the
equation for $\psi$ is solved), or describe the prescribed time evolution
of $\psi$ and all related quantities. 
Usually, it is more convenient to prescribe the plasma current density
$j_\pll$ rather than the poloidal flux $\psi.$

\vskip 1ex\noindent
\hspace*{60pt}
\begin{tabular}{l}
\begin{picture}(110,25)
\thicklines
\multiput(15,15)(0,-30){2}{\line(1,0){30}}
\multiput(0,0)(45,-15){2}{\line(1,1){15}}
\multiput(0,0)(45,15){2}{\line(1,-1){15}}
\put(25,-5){{\tt CU}}
\thinlines
\put(67,0){\vector(1,0){40}}\put(80,5){yes}
\put(30,-20){\vector(0,-1){30}}\put(15,-35){no}
\put(120,0){$j_\pll={\tt CU}\rightarrow{\tt MU}\rightarrow{\tt FP}
	\rightarrow{\tt UPL}$}
\put(120,-20){Additionally, the condition \rf{icj} is applied.}
\end{picture}
\\
\begin{picture}(110,65)
\thicklines
\multiput(15,15)(0,-30){2}{\line(1,0){30}}
\multiput(0,0)(45,-15){2}{\line(1,1){15}}
\multiput(0,0)(45,15){2}{\line(1,-1){15}}
\put(25,-5){{\tt MU}}
\thinlines
\put(70,0){\vector(1,0){40}}\put(80,5){yes}
\put(30,-20){\vector(0,-1){30}}\put(15,-35){no}
\put(120,0){$\mu={\tt MU} \rightarrow {\tt CU} \rightarrow {\tt FP} 
	\rightarrow{\tt UPL}$}
\put(120,-20){The condition \rf{icj} is not applied.}
\end{picture}
\\[50pt]
\hspace*{-60pt}The current distribution is determined \\
\hspace*{-60pt}by the steady-state condition \rf{icU}.
\end{tabular}\\[15pt]
After one of the quantities $j_\pll$ or $\mu$ is defined, all other
quantities in the list $j_\pll,\mu,\psi,U_{\rm pl}$ are calculated
using Eqs.~\rf{Ipl}, \rf{jlon} and \req{Ohm}. 
Note that, in case when an assignment of \twd{MU} is active, no
normalization is applied although $\mu(\rho_B)$ is consistent with the
total plasma current.
That is, in general, $\mu(\rho)$ has a jump (surface current) at the
plasma edge. 
In all other cases, the additional normalization condition \rf{icj} is
applied so that an assignment of \twd{CU} is valid within to a
multiplicative constant. 

\subsection{Boundary conditions}

\subsubsection{Boundary conditions for $n_e,\;T_e,\;T_i,\;f_j$}

Besides the initial condition, the boundary values 
$n_e(\rho_B,t),\;T_e(\rho_B,t),\;T_i(\rho_B,t),\;f_j(\rho_B,t)$ 
must be prescribed as functions of time.
Electron density at the plasma edge $n_e(\rho=\rho_B,t)$ is defined
according to

\begin{tabular}{p{110pt}	p{130pt}	p{151pt}}
\begin{picture}(110,35)
\thicklines
\multiput(15,15)(0,-30){2}{\line(1,0){30}}
\multiput(0,0)(45,-15){2}{\line(1,1){15}}
\multiput(0,0)(45,15){2}{\line(1,-1){15}}
\put(20,-5){{\tt NEB}}
\thinlines
\put(67,0){\vector(1,0){40}}\put(80,5){yes}
\put(30,-20){\vector(0,-1){35}}\put(15,-40){no}
\end{picture}
	&Prescribed edge density:\vspace*{-10pt}
	$$n_e(\rho_B,t)={\tt NEB}(t)$$\\[5pt]
\end{tabular} 

\begin{tabular}{p{110pt}	p{130pt}	p{163pt}}
\begin{picture}(110,25)
\thicklines
\multiput(15,15)(0,-30){2}{\line(1,0){30}}
\multiput(0,0)(45,-15){2}{\line(1,1){15}}
\multiput(0,0)(45,15){2}{\line(1,-1){15}}
\thinlines
\put(18,-5){{\tt QNB}}
\put(70,0){\vector(1,0){40}}\put(80,5){yes}
\put(30,-20){\vector(0,-1){70}}\put(15,-55){no}
\end{picture}
	&
\begin{picture}(110,25)
\thicklines
\multiput(15,15)(0,-30){2}{\line(1,0){30}}
\multiput(0,0)(45,-15){2}{\line(1,1){15}}
\multiput(0,0)(45,15){2}{\line(1,-1){15}}
\put(18,-5){{\tt QNNB}}
\thinlines
\put(70,0){\vector(1,0){40}}\put(80,5){yes}
\put(30,-20){\vector(0,-1){30}}\put(15,-35){no}
\end{picture}
&\hspace*{-6pt}Prescribed electron flux: 
\vspace*{-10pt}
$$\hspace*{-32pt}\Gamma_e(\rho_B,t)={\tt QNB}(t)+{\tt QNNB}(t)\times
n_e(\rho_B,t)$$ 
\vspace*{5pt}\\
&\hspace*{-20pt}Prescribed electron flux: 
\vspace*{-10pt}
$$\hspace*{-40pt}\Gamma_e(\rho_B,t)={\tt QNB}(t)$$\\
\end{tabular} 
\\[-30pt] \indent
\begin{tabular}{l l}
\begin{picture}(110,35)
\thicklines
\multiput(15,15)(0,-30){2}{\line(1,0){30}}
\multiput(0,0)(45,-15){2}{\line(1,1){15}}
\multiput(0,0)(45,15){2}{\line(1,-1){15}}
\put(18,-5){{\tt QNNB}}
\thinlines
%\put(30,40){\vector(0,-1){20}}%\put(15,30){no}
\put(70,0){\vector(1,0){40}}\put(80,5){yes}
\put(30,-20){\vector(0,-1){35}}\put(15,-40){no}
\end{picture}
&%\multicolumn{2}{l}
{$\hspace*{-5pt}\Gamma_e(\rho_B,t)={\tt QNNB}(t)\times 
n_e(\rho_B,t)$}\\[60pt]
\multicolumn{2}{l}{\hspace*{-20pt}Edge density is determined by}\\
\multicolumn{2}{l}{\hspace*{-20pt}a data file as
			$n_e(\rho_B,t)={\tt NEXB}(t_0)$}
\end{tabular} 
\\[20pt]
It is seen that giving ${\tt NEB=1}$ one obtains the boundary condition
$n_e(\rho_B,t)=10^{19}{\rm m}^{-3}.$
In such a case, the setting of ${\tt QNB}$ or ${\tt QNNB}$ does not
influence $n_e(\rho_B,t)$.
Otherwise, if the particle flux, rather than the particle density, at
the plasma edge is to be prescribed, then any instructions of the type 
${\tt NEB}=\cdots$ must not be present. 

Any combination of explicit or implicit approximation for the edge
particle flux can be used.
For instance, both instructions
${\tt QNB=10*NE}$ and ${\tt QNNB=10}$ describe
the same boundary condition (to within the numerical accuracy).
However, in the first case, the numerical approximation reads 
$\Gamma_e(\rho_B,t)=10n_e(\rho_B,t-\Delta t),$ while in the second,
$\Gamma_e(\rho_B,t)=10n_e(\rho_B,t).$ 

Other boundary conditions can be defined in a parallel fashion, as
illustrated in Table~4.8.
In addition to this table, we note that if none of the control parameters
listed there 
\linebreak\indent
{\bf Table 4.8. Summary of boundary condition setting}
\bigskip\\ \noindent
\begin{tabular}{| c | l | l | l | l |}\hline
Quantity/Flux & Code name  
	& \multicolumn{3}{c|}{Control parameters ~/~ Units}\\	\hline
$n_e$ / $\Gamma_e$
	&{\tt NE} / {\tt QN} 
		&{\tt NEB} ~/~ $10^{19}{\rm m}^{-3}$
			&{\tt QNB~} ~/~ $10^{19}{\rm s}^{-1}$
				&{\tt QNNB~} ~/~ ${\rm m}^3$/s	\\
$T_e$ / $q_e$	& {\tt TE} / {\tt QE}	
	&{\tt TEB} ~/~ keV & {\tt QEB~} ~/~ MJ	&{\tt QETB~} ~/~ MJ/keV\\
$T_i$ / $q_e$	& {\tt TI} / {\tt QI}
	&{\tt TIB} ~/~ keV & {\tt QIB~} ~/~ MJ	&{\tt QITB~} ~/~ MJ/keV\\
$f_1$ / $\Gamma_1$&{\tt F1} / {\tt QF1}	&{\tt F1B} ~/~ $[f_1]$
		 &{\tt QF1B} ~/~ $[f_1]\;{\rm m}^3$/s
			&{\tt QFF1B} ~/~ ${\rm m}^3$/s		\\
$f_2$ / $\Gamma_2$&{\tt F2} / {\tt QF2}	&{\tt F2B} ~/~ $[f_2]$
		 &{\tt QF2B} ~/~ $[f_2]\;{\rm m}^3$/s
			&{\tt QFF2B} ~/~ ${\rm m}^3$/s		\\
$f_3$ / $\Gamma_3$&{\tt F3} / {\tt QF3}	&{\tt F3B} ~/~ $[f_3]$
		 &{\tt QF3B} ~/~ $[f_3]\;{\rm m}^3$/s
			&{\tt QFF3B} ~/~ ${\rm m}^3$/s		\\
\hline
\end{tabular}
\\[2ex]
appears in a model, then a boundary value from the data file is used.
However, there is a difference between the following two cases: (i)
where the assignment ${\tt NEB=NEXB}$ is explicitly given in a model,
and (ii) where the same assignment is used because none of the control
parameters appears in a model. 
In the first case, the assignment is considered as time dependent while
in the second case it is executed only at the initial time.
Finally, if a data file does not define the quantity either, then the
boundary value is set to one of the defaults given in Table~4.7.

\subsubsection{Boundary condition for the poloidal flux $\psi=\tt FP$}

The boundary condition for the poloidal flux is defined according to
the following diagram

\begin{tabular}{p{110pt}	p{130pt}}
\begin{picture}(110,25)
\thicklines
\multiput(15,15)(0,-30){2}{\line(1,0){30}}
\multiput(0,0)(45,-15){2}{\line(1,1){15}}
\multiput(0,0)(45,15){2}{\line(1,-1){15}}
\put(20,-5){{\tt IPL}}
\thinlines
\put(67,0){\vector(1,0){40}}\put(80,5){yes}
\put(30,-20){\vector(0,-1){45}}\put(15,-40){no}
\end{picture}
&
Prescribed current:
$$\left.\prti{\psi}{\rho}\right|_{\rho=\rho_B}
=\frac{0.4\pi}{G_2}{\tt IPL}(t)$$
\end{tabular}
%\vspace*{-10pt}

\begin{tabular}{p{110pt}	p{124pt}	p{169pt}}
\begin{picture}(110,25)
\thicklines
\multiput(15,15)(0,-30){2}{\line(1,0){30}}
\multiput(0,0)(45,-15){2}{\line(1,1){15}}
\multiput(0,0)(45,15){2}{\line(1,-1){15}}
\thinlines
\put(18,-5){{\tt LEXT}}
\put(70,0){\vector(1,0){40}}\put(80,5){yes}
\put(30,-20){\vector(0,-1){90}}\put(15,-65){no}
\end{picture}
&
\begin{picture}(110,25)
\thicklines
\multiput(15,15)(0,-30){2}{\line(1,0){30}}
\multiput(0,0)(45,-15){2}{\line(1,1){15}}
\multiput(0,0)(45,15){2}{\line(1,-1){15}}
\put(18,-5){{\tt UEXT}}
\thinlines
\put(70,0){\vector(1,0){40}}\put(80,5){yes}
\put(30,-20){\vector(0,-1){30}}\put(15,-35){no}
\end{picture}
&\hspace*{-15pt}External circuit equation is solved: 
\vspace*{-10pt}
$${d\over dt}\left(L_{\rm ext}I\right)+U_{\rm plB}={\tt UEXT}(t)$$
\vspace*{5pt}
\\
&\vspace*{-10pt}\hspace*{-20pt}Short-circuited primary coil: 
\vspace*{-10pt}
$$\hspace*{-20pt} {d\over dt}\left(L_{\rm ext}I\right)+U_{\rm plB}=0$$
%$$\hspace*{-20pt}{d\over dt}\left({\tt LEXT*IPL}\right)+{\tt UPLB}=0$$
\end{tabular}

%\vspace*{-10pt}
\begin{tabular}{p{110pt}	p{140pt}}
\begin{picture}(110,25)
\thicklines
\multiput(15,15)(0,-30){2}{\line(1,0){30}}
\multiput(0,0)(45,-15){2}{\line(1,1){15}}
\multiput(0,0)(45,15){2}{\line(1,-1){15}}
\put(18,-5){{\tt UEXT}}
\thinlines
\put(70,0){\vector(1,0){40}}\put(80,5){yes}
\put(30,-20){\vector(0,-1){35}}\put(15,-40){no}
\end{picture}
&Prescribed loop voltage:\vspace*{-10pt} 
$$U_{\rm plB}=\left.\prti{\psi}{t}\right|_{\rho=\rho_B}={\tt UEXT}(t)$$
\\
\multicolumn{2}{l}{\hspace*{-20pt}The plasma current $I_{\rm pl}$ is}\\
\multicolumn{2}{l}{\hspace*{-20pt}prescribed by a data file:}\\
\vspace*{-20pt}
$$\hspace*{-20pt}\left.\prti{\psi}{\rho}\right|_{\rho=\rho_B}
=\frac{0.4\pi}{G_2}{\tt IPLX}$$\\
\end{tabular}
\\
In the simplified external circuit equation in this diagram 
$L_{ext}=\tt LEXT$ can be understood as the external inductance between
plasma and primary coil. 
Note also that in the case of prescribed loop voltage the parameter
{\tt UEXT} defines a loop voltage at the plasma surface and coincides
with $U_{\rm plB}=\tt UPLB$, whereas in the circuit equation {\tt UEXT}
denotes a voltage provided by the external source (primary coil). 
In the code, the circuit equation is represented as
$$
%U_{\rm ext}=\prti{\psi}{t}
%+{1\over 0.4\pi}\prti{}{t}\left(L_{\rm ext}G_2\prti{\psi}{\rho}\right)
%,\qquad
\lb.{L_{\rm ext}G_2\over 0.4\pi}\:\prti{\psi}{\rho}\rb|_{\rho_B}
+\lb.\psi\rb|_{\rho_B}=\int\limits_0^tU_{\rm ext}dt
$$
with $L_{\rm ext}$ measured in $\!\!\mu\!\!$Hn. 

\subsection{Link between the transport and equilibrium parts of Astra}

Two versions of the equilibrium solver are presently provided by the
Astra system. 
The more simple version \cite{Z3m} solves the equilibrium equation
\req{GSp} using 3-moment approach. It is assumed that the plasma
configuration is up-down symmetric and each magnetic surface can be
parameterized as 
\begin{equation}\label{3mom}
\lb\{\begin{array}{l}
r(a,\theta)=R_0+\Delta(a)+a\lb(\cos\theta-\delta(a)\sin^2\theta\rb)\\
z(a,\theta)=\Delta_Z+a\lambda(a)\sin\theta
\end{array}\rb.
\end{equation}
where $a$ is a minor radius of every magnetic surface in the equatorial
mid-plane. 
From now on $a$ will denote this magnetic surface function. 
$\Delta$ is the Shafranov shift, $\lambda$ is the elongation and
$\delta$ is the triangularity of each flux surface. 
The parameter $\Delta_Z$ is not included in the equilibrium solver
because the problem has translational symmetry in the $z$--direction. 
However, the plasma up-- or down-shift may be relevant for the
calculation of heat and particle sources.
In this representation the
boundary surface  $\mbox{\msym S}_B$ is specified with five parameters:
$R_0$, $a_B$, $\Delta(a_B)$, $\lambda(a_B)$, $\delta(a_B)$. 
This is equivalent to a prescription of four characteristic points on
the boundary. 
These are two points in the mid--plane 
$\{r=R_0+\Delta(a_B)\pm a_B,\;z=\Delta_Z\}$ 
and the two points most remote from the plane
$\{r=R_0+\Delta(a_B)-a_B\delta(a_B),\;z=\Delta_Z\pm a_B\lambda(a_B)\}$.

A more advanced equilibrium solver called
ESC\footnote{http://w3.pppl.gov/topdac/esc.htm}  
(Equilibrium and Stability Code) is also available \cite{ESC} 
which allows for arbitrary shape of the plasma boundary. 
In this version, the plasma boundary $\mbox{\msym S}_B$ is specified by
a set of $N_B\leq 121$ points $\{r_i,z_i\}$ 
while the configuration is sought in the form
\vspace*{-5pt}
\begin{equation}\label{sesconf}
\lb\{\begin{array}{l}
r(\rho,\tau)
=r_0^c(\rho)+2\ds\sum\limits_{k=1}^{K}
\lb[r_k^c(\rho)\cos k\tau+r_k^s(\rho)\sin k\tau\rb]
,\medskip\\
z(\rho,\tau)=\Delta_Z(\rho)+\rho\lambda(\rho)\sin\tau
.\end{array}\rb.
\end{equation}
ESC also provides an analysis of plasma stability with respect to tearing
and ballooning modes. 
%By request, it also creates an input file for the specialized 
%stability analysis codes PEST and DCON. 

Now we will describe the interface between the equilibrium and the
transport packages of the code.  
If the number $N_B$ does not exceed 4 and the boundary is up down
symmetric then it is convenient to use shift, elongation and
triangularity of the boundary. 
Two radially dependent functions on the right hand side of the
Grad-Shafranov equation, namely the pressure and current
density distributions, are supplied at a $\rho_j$-grid of $N_E$ points
(variable name {\tt NEQUIL}) which is different from the transport
grid. 
In the current version, $N_E={\tt NEQUIL}\leq 41$ and in addition to
its main meaning it is also used as a control parameter:
$$
\begin{array}{rll}
&N_E<0
	&\mbox{No equilibrium solver is called. Pre-calculated metric}\\
	&&\mbox{is taken from a data file as described in Section~5.3.2.}
\medskip\\
&N_E=0
	&\mbox{No equilibrium solver is called. Parabolic distributions}\\
	&&\mbox{for }\Delta(\rho),\lambda(\rho),\delta(\rho)
	\mbox{ are used and }a(\rho)=\rho/\sqrt{\lambda}.
\medskip\\
0 < &N_E\leq 41
	&\mbox{3-moment equilibrium solver is called with $N_E$-point grid.}
\medskip\\
&N_E>41\quad
	&\mbox{ESC is called.}
\end{array}
$$
%The information in the table above is sufficient for the work with the
%Astra code if the 4-point set prescribes the plasma boundary. 
%More knowledge about the interface is needed only for more complicated
%boundary shape. 
Parameter exchange is shown in the following diagram and explained in
the table below. 
\vspace*{-25pt}
$$
\begin{array}{lcr}
\bigskip\\
\fboxrule=1pt\fboxsep=5mm
\hspace*{-5pt}\fbox{\Large\bf Transport}
&
\begin{array}{c}
R_0B_0,\;p,\;j_\pll,\;N_E,\;N_B,\;\{r_i,z_i\}\\
\vector( 1,0){172}\\[-10pt]
\vector(-1,0){172}\\
\rho_B,\;\rho_j({\bf r}),\;I,\;V',\;G_1,\;G_2,\;G_3\\
\mbox{\rm Additional output is available}\\
\mbox{\rm by request}
\end{array}
&
\fboxrule=1pt\fboxsep=5mm
\fbox{\Large\bf Equilibrium}
\bigskip\\
\end{array}
$$

%$$
%\begin{array}{l}
%\hspace*{0.6em}
\noindent
\begin{tabular}{|ll|ll|}\hline &&&\\
\multicolumn{2}{|c}
	{\hspace*{1.5em}{\bf Input} to the equilibrium module}\hspace*{1.5em}
&\multicolumn{2}{|c|}
	{\hspace*{1.5em}{\bf Output} from the equilibrium module 
	 \hspace*{1em}}\\[10pt]\hline
$R_0B_0$	& boundary value of $I$
	&$\rho_B$		& boundary value of $\rho$\\[8pt]

$N_B$		& boundary-grid size ($N_B\leq 12$)
	&$\rho_j({\bf r})$	& shape of every flux surface\\[8pt]

$\{r_i,z_i\}$	& boundary shape ($1\leq i\leq N_B$)
	&$I(\rho_j)$		& poloidal current \\[8pt]

$N_E$		& $\rho$-grid size and type of solver
	&$V'(\rho_j)$		& volume derivative\\[8pt]

\hspace*{-0.5em}
	\begin{tabular}{l}
	$\rho_j$\\[8pt]
	$p(\rho_j)$\\[8pt]
	$j_\pll(\rho_j)$\\[8pt]
	\end{tabular}
	    &	\hspace*{-0.5em}
		\begin{tabular}{l}
		$\rho$-grid ($1\leq j\leq N_E$)\\[8pt]
		pressure profile\\[8pt]
		current density profile\\[8pt]
		\end{tabular}
		    &	\hspace*{-1em}
			$\lb.\begin{tabular}{l}
			$G_1(\rho_j)$\\[8pt]
			$G_2(\rho_j)$\\[8pt]
			$G_3(\rho_j)$
			\end{tabular}\rb\}$
			   &	\begin{tabular}{l}\hspace*{-0.5em}
				metric characteristics \\
				\end{tabular}\\
\hline
\end{tabular}%\bigskip\\
%\end{array}
%$$

\subsection{Astra variables}

In this section, we describe Astra internal variables dividing them into
groups according to their meaning and limiting the discussion to those
variables which can be accessed (or corrupted) by a user of the
code. 
These variables are held in Fortran common blocks and seen from every
part of the code. 
First of all, there are variables describing different physical
quantities: some of them have to be defined by a user, others are
calculated by the code. 
The second group consists of control parameters for the numerical
schemes used in the code and for tuning the output.
Finally, there is a set of global variables without any predefined
meaning, which are entirely at the disposal of the user. 
Employing these variables is recommended if any new quantity
is needed in a simulation. 
Although it is not forbidden to introduce new variables in a user's
model, no check is made for dangerous possible name conflicts. 
Therefore, it is supposed that in introducing a new variable (see
Section~5.2.8), the user will check that the variable name does not
coincide with any of the Astra global variables.
Complete lists of the Astra global variables can be found in the files:
\footnote
{We remind that \twd{AWD} stands here for the Astra working directory
which can be different for each installation.}\label{AWD} 
%
\begin{tabbing}
\qquad\tw{AWD/for/parameter.inc},\qquad\=\tw{AWD/for/const.inc},
\qquad\=\qquad\tw{AWD/for/status.inc}, \kill
\qquad\tw{AWD/for/parameter.inc},\>	\tw{AWD/for/const.inc}, \>
\tw{AWD/for/status.inc}, \\
\qquad\tw{AWD/for/outcmn.inc},\>	\tw{AWD/tmp/declar.fml},\>
\tw{AWD/tmp/declar.fnc}. 
\end{tabbing}

It is obvious that there is a minimal set of variables which must be
defined by a user in order to start an Astra simulation. 
In the Astra code we have tried to minimize this burden on the user by
setting as many variables as possible according to common sense.
We will use the following symbols to mark different types of the code
variables.\bigskip \\ 
\begin{tabular}{c p{151mm}}
  !  & The variable must be defined by a user. 		\medskip \\
  +  & The variable can be defined by a user. Otherwise, it takes
	a default value. 				\medskip \\
  $*$ & The variable cannot be defined by a user. It is always defined 
	by the code. 					\bigskip \\
%  -- & The variable cannot be defined by a user. It is defined 
%	by the code if needed. 				\medskip \\
%$\uplus$& The variable is reserved for use in some special subroutines. 
%	However, it can be used for other goals if the corresponding 
%	subroutine is not called.			\bigskip\\
%\tt * & The variable can be used arbitrary.
\end{tabular}\\
As mentioned, there is an additional type of variables which are
exclusively defined by the user and can be used for exchanging
information between user routines.

There are several ways for determining variables in the Astra
code and the default definition is only applied if the variable
was not explicitly set by the user. 
More about possible ways of setting code variables and about their
default values can be found in Section~5.3.

\subsubsection{Scalar (time dependent) variables}

We start with a set of variables which define the plasma configuration.
The five geometrical variables \twd{AB, ELONM, TRICH} and \twd{AWALL,
RTOR} are time independent.
The first three of them are the maximum possible values for the time
dependent quantities \twd{ABC, ELONG} and \twd{TRIANG}, respectively. 
The variable \twd{AWALL} is  presently used only for a radial grid
allocation 
\linebreak
%\begin{figure}
\indent
{\bf Table 4.9. Configuration parameters}\\[2ex]
\begin{tabular}{|c|l|c|l|}\hline
\hspace*{-1pt}\bf Type\hspace*{-1pt} & \bf Variable 
			&\bf Units &\bf Description\\		\hline
 ! &\tt AB    & m & Limiter position in the mid-plane, $a_B$\\
 + &\tt ABC   & m & Transport grid boundary in the variable $a$\\
 + &\tt AWALL & m & Wall position in the mid-plane\\
 ! &\tt RTOR  & m & Major radius of the vacuum chamber\\
 + &\tt SHIFT & m & Horizontal shift of the edge magnetic surface\\
 + &\tt UPDWN & m & Vertical shift of the edge magnetic surface\\
 + &\tt ELONG & -- & Elongation of the edge magnetic surface\\
 + &\tt ELONM & -- & Vacuum chamber elongation (used for drawing)\\
 + &\tt TRIAN & -- & Triangularity of the boundary surface\\
 + &\tt TRICH & -- & Vacuum chamber triangularity	\\	\hline
\end{tabular}
%\end{figure}
\bigskip\\
%\begin{figure}
\indent{\bf Table 4.10. Main plasma parameters}\\[2ex]
\begin{tabular}{| c | l | c | l |}\hline
\hspace*{-1pt}\bf Type\hspace*{-1pt} & \bf Variable 
		& \bf Units & \bf Description \\	\hline
 ! &\tt BTOR &  T  & Vacuum toroidal field at the position $\tt RTOR$\\
 ! &\tt IPL  &  MA & Plasma current\\			
 + &\tt AMJ  &$m_p$& Main ion atom mass in units of the proton mass \\
 + &\tt ZMJ  &$e_p$& Main ion charge in units of the proton charge\\
\hline
\end{tabular}\bigskip\\%\vskip 3ex
%\end{figure}
and usually it does not require an explicit definition.
The major radius \twd{RTOR} corresponds to 
the position where the vacuum toroidal magnetic \twd{BTOR} is given. 
Although \twd{RTOR} cannot vary 
in time, plasma motion can be defined by varying \twd{SHIFT} and
\twd{UPDWN}.~ 
The main ion mass and charge can also be given as radial arrays
\twd{AMAIN} and \twd{ZMAIN}.

{\bf Table 4.11. Other scalar plasma parameters}\\[2ex]
\begin{tabular}{| c |l|c|l|}\hline
\hspace*{-1pt}\bf Type\hspace*{-1pt} & \bf Variable 
			&\bf Units &\bf Description	\\	\hline
 + &\tt GN2E &	-- & Factor in convective electron heat flux\\
 + &\tt GN2I &	-- & Factor in convective ion heat flux\\
 + &\tt AIM1,AIM2,AIM3
	  & $m_p$ & Impurity species mass in units of the proton mass\\
 + &\tt LEXT &$\mu$H & External inductance\\
 + &\tt UEXT &   V   & External loop voltage\\
%\tt QETB  & MJ/keV & Coefficient for electron heat outflux = QETB*TEB\\
%\tt QFF1B &  1/s   & Coefficient for F1 outflux = QFF1B*F1B\\
%\tt QFF2B &  1/s   & Coefficient for F2 outflux = QFF2B*F2B\\
%\tt QFF3B &  1/s   & Coefficient for F3 outflux = QFF3B*F3B\\
%\tt QITB  & MJ/keV & Coefficient for ion heat outflux = QITB*TIB\\
%\tt QNNB  & $m^3/s$& Coefficient for density outflux = QNNB*NEB\\
 + &\tt WNE, WTE, WTI& m & Width of exponential decay 
		for $n_e,T_e,T_i$ at $a>a_B$\\
\hline
 $*$ &\tt RHOW   & m & 
    Limiter position in the variable $\rho$: $a(\tw{RHOW})=\tw{AB}$\\
%\tt ROWALL&   m  & Wall position    (in variable $\rho$)\\
 $*$ &\tt ROC   & m  & $\rho_B,$
    grid edge in the variable $\rho$: $a(\tw{ROC})=\tw{ABC}$\\
 $*$ &\tt VOLUME &m$^3$ & Plasma volume inside $\rho_B=$\tw{ROC}\\
 $*$ &\tt VSB   & Vs & Poloidal flux at the plasma edge \\
 $*$ &\tt VSC   & Vs & Poloidal flux at the magnetic axis\\	\hline
\end{tabular}\vskip 3ex

As seen from Table~4.11, up to three different ion impurity species
can be defined in the code by their masses, charges and densities
\footnote[1]{Unlike the masses, the impurity charges and densities are
described by radial arrays {\tt~ZIM1, ZIM2, ZIM3} and
\twd{NIZ1,NIZ2,NIZ3} (Section~4.6.2).}.
The impurities are not directly included in the main set of transport
equations, but they are required in the NBI heating subroutine, NCLASS,
STRAHL and some other external modules. 
The parameters \twd{WNE, WTE, WTI} are used to describe the plasma
parameters \twd{NE, TE, TI} in the scrape-off-layer outside the main
transport grid.
In the present version, they do not play any role in the simulation and
can be seen only in a special drawing mode when radial profiles are
plotted beyond the transport boundary. 

As long as the memory requirement for the Astra code is quite moderate
the code does not use dynamic memory allocation. 
All arrays have fixed dimensions and any change of those requires
re-installation of the code. 
Some of the code constants are given in the Table~4.12.
Normally, an error is reported when any of these array limits is
exceeded. 

The parameter $\tt NRD$ allocates the maximum possible length of all
radial arrays. 
However, the parameter which really determines the radial grid size is
$\tt NB1\leq NRD.$ 
The transport core of the code is quite fast so that large values of
\twd{NB1} do not affect its speed significantly. 
\linebreak\indent
{\bf Table 4.12. Internal constants}\\[2ex]
\begin{tabular}{| c | l|c|l|}\hline
\hspace*{-1pt}\bf Type\hspace*{-1pt} & \bf Constant 
			&\bf Value &\bf Description	\\	\hline
% + &\tt SHOT  &	     & String with a shot description\\
% + &\tt DEVICE&	     & String with a device description\\
 $*$ &\tt NRD   & 501  & Maximum size of the radial grid\\
 $*$ &\tt NRDX  & 200  
	& Maximum No. of radial points for input from a data file \\
 $*$ &\tt NRW   & 96   & Maximum No. of radial / time curves for output\\
 $*$ &\tt NTIMES& 1024 
	& Dimension of arrays for output in time modes 6 and 7\\
 $*$ &\tt NTVAR & 2000 
	& Total No. of time slices for all variables in a data file\\
 $*$ &\tt NTARR & 5000 
	& Total No. of time slices for all arrays in a data file\\
 $*$ &\tt GP    & $\pi$ & \\
 $*$ &\tt GP2   & $2\pi$   & \\
\hline
\end{tabular}\\[2ex]
However, some additional modules which do not use a separate grid can
noticeably decelerate the calculations.
The essential feature of the parameter \twd{NB1} is that it must be
defined at the earliest phase of the code execution and cannot be
redefined later.
It means that, unlike most other parameters in the Astra code which can
be set in many different ways, the only way of defining the grid size
\twd{NB1} is input from a data file (Section~5.3.1).
The parameter \twd{NEQUIL} defines a radial grid for the 3-moment
equilibrium solver and serves also as a switch to the $\tt{ESC}$ solver. 
In the latter case, a radial grid is selected automatically. 
Two parameters \twd{XFLAG} and \twd{NBND} define a type of the plasma
boundary and a number of points on it, respectively.

%\linebreak\indent
{\bf Table 4.13. Parameters for radial grid control}\medskip\\
\begin{tabular}{| c |l|c|l|}\hline
\hspace*{-1pt}\bf Type\hspace*{-1pt} & \bf Variable 
			&\bf Units &\bf Description	\\	\hline
+ &\tt NB1    & -- & No. of points of the main transport grid 
					\twd{NB1}$\leq$\twd{NRD}\\
+ &\tt NEQUIL & --& 3-moment equilibrium solver grid size and control switch\\
+ &\tt NB2EQL & --& Factor for NB pressure contribution 
					to equilibrium eqn \\
+ &\tt NBND   & --& No. of points on the boundary\\
+ &\tt XFLAG  & --& X-point flag				\\
+ &\tt NUF    & --& No. of radial points to be written in a U-file\\
+ &\tt XOUT   & --& Type of abscissa for radial graphic output 
						(modes 1, 2, 4, 5)\\
+ &\tt XINPUT & --& Type of abscissa for radial profile input\\ \hline
$*$ &\tt HRO    & m & Radial grid step in the variable $\rho$\\
$*$ &\tt HROA   & m & Edge step of the radial grid: 
					\tt RHO(NA1)-RHO(NA)\\
$*$ &\tt NA     & --& = \tw{NA1}-1\\
$*$ &\tt NA1    & --& Edge grid point number: {\tt ROC=RHO(NA1)}\\
$*$ &\tt NAB    & --& $\tw{NAB}\geq\tw{NA1},\;\;\tw{AB}=\tw{AMETR(NAB)}
		\geq\tw{ABC}=\tw{AMETR(NA1)}$\\			\hline
\end{tabular}\vskip 3ex\noindent

The parameter \twd{NUF} defines a number of radial grid points for an
output in the format of U-file. 
A type of the created U-file depends on the current radial graphic
mode (Section~5.5.3)\\
\indent -- 1D (radial coordinate) in graphic modes 1, 2, 3,\\
\indent -- 2D (radial and time coordinates) in the graphic mode 4,\\
\indent -- 1D (time coordinate) in the graphic mode 6.\\
In addition, one should also have in mind that the radial coordinate in
the U-file produced, is the same as for the current graphs on the
screen. 
Therefore, it is controlled by parameter \twd{XOUT} which
defines an abscissa type for radial profiles in the current plot. 
If the parameter \twd{XOUT} takes values 0 or 1 (default value is 1)
then the flux surface label $a$ is used as a radial coordinate. 
If $\tw{XOUT}=2$ or 3 then $\rho$ or $\psi$ are used in place of $a.$
The parameter \twd{XINPUT} can take the same values, but it defines an
abscissa of radial profiles for input in cases when the radial grid is
equidistant but the coordinate is not defined explicitly.
The meaning of these numbers is similar to that for the parameter
\twd{GRIDTYPE} 
which is explained in more detail in Section~5.4.4, and summarized in
Table~5.5. 
%Finally, we note that for a 2D U-file output the only radial coordinate,
%$0<a<\tt AB,$ is meaningful. 
All other parameters in Table~4.13 cannot be changed by a user.
They may be used for output only.

{\bf Table 4.14. Parameters for output control}\\[10 pt]
\begin{tabular}{|c|l|c|l|}\hline
\hspace*{-1pt}\bf Type\hspace*{-1pt} & \bf Variable 
	&\bf\hspace*{-1pt}Units\hspace*{-1pt} &\bf Description\\ \hline
 + &\tt DPOUT  & s 
   & Time interval for profile storing (modes 4, 5 and Review)\\
 + &\tt DROUT  & s 
   & Time interval for profile redrawing (modes 1, 2 and 3)\\
 + &\tt DTOUT  & s & Time step for time curve output (modes 6 and 7)\\
 + &\tt TPAUSE & s  & Time for switching on the ``step'' mode	\\
 + &\tt TEND   & s  & End of run time				\\
 + &\tt TINIT  & s & Time axis left label (modes 6 and 7)   \\
 + &\tt TSCALE & s & Time axis scale length (modes 6 and 7) \\	\hline
\end{tabular}\vskip 3ex

Table~4.14 contains a number of other control parameters for graphic
output. 
All these parameters are described in Section 5.5.3.
They can be set in a model or adjusted interactively.
If the parameters \twd{DPOUT}, \twd{DROUT}, \twd{DTOUT}, \twd{TSCALE} and
also the parameters \twd{TAUMIN, TAUMAX} from the Table~4.15 are not
given by the user they are set by the code to be proportional to the
plasma volume.

The control parameters given in Table~4.15 define the numerical algorithm
used in the transport part of the code. 
Usually, the time step is selected automatically by the following
condition. 
If a maximum relative change in any of variables \tw{~TE, TI, NE, F1,
F2, F3~} is larger than \twd{DELVAR} then the time step $\tau=\tt TAU$
is decreased. 
Otherwise, it is multiplied by \twd{TAUINC} which when not defined
explicitly is equal to 1.1. 
However, a user can fix the time step by setting ${\tt TAUMIN} = {\tt
TAUMAX},$ or modify the algorithm (see Section~4.9.8). 
An increase of the parameter $\tt ITEREX$ (default value 1) can
sometimes improve the numerical stability in strongly nonlinear problems.
Another iteration parameter $\tt ITERIN$ cannot be changed by a user.
It is selected by the code to ensure that the numerical approximation of
electron--ion heat exchange does not violate energy conservation. 

{\bf Table 4.15. Parameters for accuracy control}\\[2ex]
\begin{tabular}{| c |l|c|l|}\hline
\hspace*{-1pt}\bf Type\hspace*{-1pt} & \bf Variable 
			&\bf Units &\bf Description		\\ \hline
 + &\tt TAUMAX & s  & Maximum time step: $\tau\leq{\tt TAUMAX}$	\\
 + &\tt TAUMIN & s  & Minimum time step: $\tau\geq{\tt TAUMIN}$	\\
 + &\tt TAUINC & --& Maximum time step increment: 
		${\tt TAU}(t+\tau)/{\tt TAU}(t)$		\\
 + &\tt DELVAR & --& Maximum relative change of variables at time step\\
 + &\tt ITEREX & --& No. of iterations in the external loop 	\\ \hline
$*$&\tt ITERIN& --& No. of iterations in the internal loop 	\\
$*$&\tt TIME  & s  & Current time of the simulation		\\
$*$&\tt TAU   & s  & Current time step				\\
% + &\tt DTEQ  & s  & Subroutine J call interval     DTEQ(1,J)	\\
% + &\tt TEQ   & s  & Subroutine J last call time    TEQ(J)	\\ 
\hline
\end{tabular}\vskip 3ex

\paragraph{Exchange variables.}
A group of parameters, which are in all respects similar to those
discussed above, is reserved for particular uses in the different 
external modules rather than in the Astra core.
If, in a particular model, these modules are not involved then the
parameters can be used quite arbitrarily independently of their names and
significance.
If the modules are in use, however, then the parameters can only be used
in the predefined context. 

\indent
{\bf Table 4.16. Input/output parameters for the subroutine \twd{NEUT}}
\\[2ex]
\begin{tabular}{|l|c|l|}					\hline
\bf Name		&\bf Units &\bf Description	\\	\hline
\tt NNCL / NNWM &$10^{19}/{\rm m}^{-3}$ 
	& Density of incoming cold / warm neutrals\\
\tt ENCL / ENWM & keV 
	& Energy of incoming cold / warm neutrals\\
\tt NNCX  & -- & No. of iterations in the neutral solver\\
\tt ALBPL& -- 
	& Plasma albedo (Neutral flux outward)/(Neutral flux inward)\\
\hline
\end{tabular}%\bigskip\\
\vskip 3ex

The scalar input parameters for the external module 
$\tt NEUT$ are listed in Table~4.16.
Two pairs $\{{\tt NNCL, ENCL}\}$ and ${\tt \{NNWM,ENWM\}}$ are available.
Each pair describes a density and energy of one mono-energetic incident
neutral component.
Although it is not quite correct, the parameter \twd{NNCX} can be
viewed as a number of charge exchange neutral 
generations.
The complete list 
of parameters related to this subroutine which
includes also output is discussed in Section~4.11. 

A number of control parameters are used in different heating and current
drive modules. 
Most of them are built into the appropriate packages, and can be
changed by the user according to the usual rules. 
Only a few of them are associated with the total heating power  
\twd{(QECR, QFW, QICR, QLH, QNBI)} and these
belong to the list of main Astra parameters.
As already mentioned, the meaning of these parameters is defined
inside the corresponding packages and can even change if one heating
module is replaced by another. 

All scalar variables discussed so far can be defined with maximum
flexibility (except for those marked with `$\!\!*\!\!$'). 
They can be defined by all the methods described in Section 5.3, i.e. in
a model, in a data file (with few exceptions) and interactively. 
Other scalars described below have fewer restrictions in 
their usage, but have less flexibility in how they are initialized.

\paragraph{Dummy variables.}
As mentioned in the beginning of this section there are a number of
parameters which are reserved for user's needs and can be used 
arbitrarily in a model. 
All these scalar variables have no predefined meanings and mnemonics of
names (if any) should not be viewed as binding.
For instance, \twd{CMHD1} need not be associated with MHD phenomena as
well as \twd{CIMP1} may have nothing to do with the impurity behavior.
They can be used for communication between different user's
subroutines, in intermediate calculations or for output. 
Two different groups are described below. 
Variables in the both groups can be time dependent but the time
dependence is to be set in different ways.

\indent
{\bf Table 4.17. Parameters for a run control (C-parameters)}\\[2ex]
{\tt\begin{tabular}
{|p{15mm}p{15mm}p{15mm}p{15mm}p{15mm}p{15mm}p{15mm}p{15mm}|}       \hline
CF1	& CF2	& CF3	& CF4	& CF5	& CF6	& CF7	& CF8	\\
CF9	& CF10	& CF11	& CF12	& CF13	& CF14	& CF15	& CF16	\\
CV1	& CV2	& CV3	& CV4	& CV5	& CV6	& CV7	& CV8	\\
CV9	& CV10	& CV11	& CV12	& CV13	& CV14	& CV15	& CV16	\\
CHE1	& CHE2	& CHE3	& CHE4	& CHI1	& CHI2	& CHI3	& CHI4	\\
CNB1	& CNB2	& CNB3	& CNB4	& CNBI1	& CNBI2	& CNBI3	& CNBI4	\\
CCD1	& CCD2	& CCD3	& CCD4	& CRF1	& CRF2	& CRF3	& CRF4	\\
CNEUT1	& CNEUT2& CNEUT3& CNEUT4& CPEL1	& CPEL2	& CPEL3	& CPEL4	\\
CBND1	& CBND2	& CBND3	& CBND4	& CFUS1	& CFUS2	& CFUS3	& CFUS4	\\
CIMP1	& CIMP2	& CIMP3	& CIMP4	& CMHD1	& CMHD2	& CMHD3	& CMHD4	\\
CRAD1	& CRAD2	& CRAD3	& CRAD4	& CSOL1	& CSOL2	& CSOL3	& CSOL4	\\ \hline
\end{tabular}}\vskip 3ex

The group of parameters listed in Table~4.17 will be called
C-parameters.
They can be used for run control because
the C-parameters can be changed interactively during a code run.  
They can also be saved during a run (Section~5.3.2), 
then the subsequent run will be started with the recently adjusted
values of these control parameters. 
If the C-parameters are not determined explicitly in one of the ways
described in Section~5.3 then they take default values which are 0 
for the 16 variables starting with \twd{CV} (from \twd{CV1} till
\twd{CV16}) and 1 for all others. 
A limitation is that the C-parameters cannot be read from a data file. 

A second set of dummy parameters given in Table~4.18 can be defined by
an input file. 
These parameters can be used to input non-standard experimental
characteristics for which 
\linebreak\indent
{\bf Table 4.18. Variables readable from a data file}\\[2ex]
{\tt\begin{tabular}
{|p{15mm}p{15mm}p{15mm}p{15mm}p{15mm}p{15mm}p{15mm}p{15mm}|}      \hline
ZRD1X & ZRD2X & ZRD3X & ZRD4X & ZRD5X & ZRD6X & ZRD7X & ZRD8X \\
ZRD9X & ZRD10X& ZRD11X& ZRD12X& ZRD13X& ZRD14X& ZRD15X& ZRD16X\\ 
ZRD17X& ZRD18X& ZRD19X& ZRD20X& ZRD21X& ZRD22X& ZRD23X& ZRD24X\\
ZRD25X& ZRD26X& ZRD27X& ZRD28X& ZRD29X& ZRD30X& ZRD31X& ZRD32X\\  
ZRD33X& ZRD34X& ZRD35X& ZRD36X& ZRD37X& ZRD38X& ZRD39X& ZRD40X \\
ZRD41X& ZRD42X& ZRD43X& ZRD44X& ZRD45X& ZRD46X& ZRD47X& ZRD48X \\
\hline
\end{tabular}}\bigskip\\
(unlike \twd{BTOR} or \twd{IPL}) no special variable is allocated. 
Typical examples are measured 
$q=2$ radius or a diamagnetic signal.
Default value for all these parameters is $0$ and they cannot be changed
interactively.
This set of parameters is coupled with the next set by the rules 
\linebreak\noindent
\indent
{\bf Table 4.19. Z-parameters}\\[2ex]
{\tt\begin{tabular}
{|p{15mm}p{15mm}p{15mm}p{15mm}p{15mm}p{15mm}p{15mm}p{15mm}|}	\hline
ZRD1	& ZRD2	& ZRD3	& ZRD4	& ZRD5	& ZRD6	& ZRD7	& ZRD8	\\
ZRD9	& ZRD10	& ZRD11	& ZRD12	& ZRD13	& ZRD14	& ZRD15	& ZRD16	\\ 
ZRD17	& ZRD18	& ZRD19	& ZRD20	& ZRD21	& ZRD22	& ZRD23	& ZRD24	\\
ZRD25	& ZRD26	& ZRD27	& ZRD28	& ZRD29	& ZRD30	& ZRD31	& ZRD32 \\ 
ZRD33 	& ZRD34 & ZRD35 & ZRD36 & ZRD37 & ZRD38 & ZRD39 & ZRD40 \\
ZRD41 	& ZRD42 & ZRD43 & ZRD44 & ZRD45 & ZRD46 & ZRD47 & ZRD48 \\ \hline
\end{tabular}}\bigskip\\
discussed in Section~5.3.3. 

\subsubsection{Vector (radially and time dependent) variables} 

A few hundreds of internal Astra variables describe radially and time
dependent quantities. 
Here we discuss those which are implemented as Fortran arrays,
i.e. those which are calculated once at each time step at all points on
a radial grid.
Other radially dependent quantities will be considered in the
next section.

A complete alphabetical list of all Astra arrays is given in the file~
\tw{AWD/for\hspace*{-.5pt}/\hspace*{-.5pt}status\hspace*{-.5pt}.\hspace*{-.5pt}inc}. 
In this section, for the sake of convenience, the arrays are combined in
groups according to their physical meaning. 
A group of vector variables related to the plasma geometry is given in
Table~4.20.
The array $\tw{RHO}=\rho_j=(j-0.5)\;\tw{HRO},$ $(j=1,2,3,\dots,\tt NA)$ 
is the main radial 
\linebreak
\indent
{\bf Table 4.20. Arrays describing plasma configuration}\bigskip\\
\begin{tabular}{| c | l | c | l  l |}\hline
\hspace*{-1pt}\bf Type\hspace*{-1pt} & \bf Variable    
		&\bf Units &\multicolumn{2}{|c|}{\bf Description}\\	\hline
$*$&\tt AMETR     & m & $a$ & Magnetic surface radius in a mid--plane\\
$*$&\tt RHO       & m & $\rho$ & Equivalent radius of a magnetic surface \\
+ &\tt SHIF / SHX& m & $\Delta$ & Shafranov shift of a magnetic surface \\
+ &\tt SHIV      & m & $\Delta_Z$ & Vertical shift of a magnetic surface\\
+ &\tt ELON / ELX&   & $\lambda$ & Elongation of a magnetic surface \\
+ &\tt TRIA / TRX&   & $\delta$ & Triangularity of a magnetic surface \\
$*$&\tt SQEPS     &   & $\sqrt{\varepsilon}$ & $\varepsilon=a/R$ 
			being the  inverse aspect ratio\\
$*$ &\tt SLAT  & ${\rm m}^2$ & $S_a$ & Toroidal surface area \\
$*$ &\tt VOLUM & ${\rm m}^3$ & $V$ & Volume within a magnetic flux surface\\
$*$ / +&\tt VR~~~/ VRX& ${\rm m}^2$ & $V'$ & on the main grid at $t$\\
$*$&\tt VRO   & ${\rm m}^2$ & $V'$ & on the main grid at $t-\tau$\\
$*$&\tt VRS   & ${\rm m}^2$ & $V'$ & on the intermediate grid at $t$\\
\hline
\end{tabular}\bigskip\\ %\vskip 3ex
grid used in the transport problem. 
The main set of unknown functions in Eqs.~\rf{main}, \rf{aux}
$n_e,T_e,T_i,\psi,f_j$ and most of other arrays are defined on this
grid.  
Fluxes and other quantities which are
derivatives of quantities defined on the main grid, for instance, the
rotational transform $\mu=\tt MU,$ are defined on the intermediate grid 
$\rho_j^{aux}=j\;\tw{HRO}.$

Some variable names for arrays have the character \twd{`X'} at the end. 
These {\tt X}-arrays can be imported from an
input file as described in Section~5.3. 
For instance, these external data can be experimentally measured radial
profiles, and can be read in for on-line comparison with the results of
simulation during the run. 
They could also be radial functions, pre-calculated by stand-alone codes 
for subsequent usage in the Astra code, e.g., heat deposition profiles
or equilibrium evolution characteristics calculated elsewhere. 
Although all \twd{X}-arrays have names associated with predefined
names for some physical quantities, most of them are independent from
their prototypes.
However, there are a few situations when some unexpected side effects are
possible. 
It usually can happen when any explicit definition is missing and the
code uses default assignment as described in Section~5.3.1.
For instance, if the transport equation \rf{Aux} for \tt F1 \rm is
involved but no initial condition for \twd{F1} is specified, then the
Astra compiler will use the default assignment \twd{F1=F1X.}
Then \twd{F1X} can be used in this sense only.
Otherwise, the variable \twd{F1X} can be used for reading arbitrary
radially and time dependent function from a data file.

In the next two tables 4.21 and 4.22, the main plasma parameters,
corresponding fluxes and sources are collected. 
Note that the quantities \twd{SNTOT, PETOT, PITOT, QN, QE, QI} are
calculated by the code. 
Any user's attempt to define them will be overridden.

\vskip 1ex\indent
{\bf Table 4.21. Arrays for the main plasma parameters}\\[2ex]
\begin{tabular}{|c|l|c|l|}					\hline
\hspace*{-1pt}\bf Type\hspace*{-1pt} & \bf Variable 
			&\bf Units &\bf Description	\\	\hline
+ / $*$ / + &\tt NE / NEO / NEX &$10^{19}{\rm m}^{-3}$ 
	& $n_e(t)$ \tw{~/~} $n_e(t-\tau)$ \tw{~/~} $n_e^{exp}(t)$ \\
+ / $*$ / + &\tt NI / NIO / NIX &$10^{19}{\rm m}^{-3}$ 
	& $n_i(t)$ \tw{~/~} $n_i(t-\tau)$ \tw{~/~} $n_i^{exp}(t)$ \\
+ / $*$ / + &\tt TE / TEO / TEX & keV & $T_e(t)$ \tw{~/~} 
	$T_e(t-\tau)$ \tw{~/~} $T_e^{exp}(t)$ 	\\
+ / $*$ / + &\tt TI / TIO / TIX & keV & $T_i(t)$ \tw{~/~} 
	$T_i(t-\tau)$ \tw{~/~} $T_i^{exp}(t)$ 	\\
+ &\tt AMAIN / ZMAIN &$m_p$\tw{~/~}$e_p$ 
		& main ion mass \tw{~/~} main ion charge\\
+ &\tt ZEF~~ / ZEFX &  & $Z_{eff}$ \tw{~/~} $Z_{eff}^{exp}$\\	\hline
\end{tabular}\vskip 2.5ex

{\bf Table 4.22. Fluxes and sources \rm (see also Table~4.6)}\\[2ex]
\begin{tabular}{|c|l|c|l|}					\hline
\hspace*{-2pt}\bf Type\hspace*{-2pt} & \bf Variable 
			& \bf Units & \bf Description\\ \hline
$*$ / +&\tt SNTOT / SNX	&$10^{19}$m$^{-3}$s$^{-1}$ 
	& Particle source $\tt SNTOT = SN+SNN*NE$	\\
$*$ / +&\tt PETOT / PEX& MW$\,{\rm m}^{-3}$ 
	& Electron heat source $\tt PETOT=PE+PET*TE$	\\
$*$ / +&\tt PITOT / PIX& MW$\,{\rm m}^{-3}$ 
	& Ion heat source $\tt PITOT=PI+PIT*TI$	\\
$*$ / +&\tt GN~~~ / GNX &$10^{19}$m$^{-2}$s$^{-1}$&$\tt GN=QN/G11,$
	\tw{GNX} is usually defined in \tw{GNXSRC}\\
$*$ / +&\tt QN~~~ / QNX   & $10^{19}\rm s^{-1}$ 
	& Electron flux through the entire magnetic surface	\\
$*$ / $*$&\tt QE~~~ / QI  & MW & Electron heat flux / Ion heat flux 	\\
\hline
\end{tabular}\vskip 2.5ex

Table~4.23 describes quantities related to the diffusion of the poloidal
field. 
The quantities \twd{FV, MV, CV} are introduced for simulation of a
stellarator. 
For tokamak modelling 
\\[1ex]\indent
{\bf Table 4.23. Poloidal flux and related quantities}\bigskip\\
\begin{tabular}{|c|l|c|l|}					\hline
\hspace*{-1pt}\bf Type\hspace*{-1pt} & \bf Variable 
			&\bf Units &\bf Description	\\	\hline
+ / $*$&\tt FP / FPO &Vs
		&Poloidal flux: $\psi(t)$ \tw{~/~} $\psi(t-\tau)$\\
+ &\tt MU / MUX & & Inverse safety factor ($\mu=\iota=1/q$)\\
+ &\tt CU / CUX & MA$\,\rm m^{-2}$& $j_\pll,$ 
		longitudinal current density, \req{dfjp}\\
+ &\tt CD   & MA$\,\rm m^{-2}$ & $j_{CD},$ driven current, \req{jCD}\\
+ &\tt CUBS & MA$\,\rm m^{-2}$ & $j_{BS},$ bootstrap current, \req{jCD}\\
+ &\tt CUTOR& MA$\,\rm m^{-2}$ & $j_{\rm tor},$ 
		toroidal current density, \req{dfjtor}\\
+ &\tt ULON &  V  & $U_\pll,$ parallel loop voltage, \req{dfEp}\\
+ &\tt UPL  &  V  & $U_{\rm pl},$ toroidal loop voltage, \req{Upl}\\
\hline
+ &\tt FV   & Vs & Vacuum poloidal flux created by external coils\\
+ &\tt MV / MVX & & Vacuum rotational transform for a stellarator (iota)\\
+ &\tt CV   &MA$\,\rm m^{-2}$& Current density equivalent to \tw{MV}\\
\hline
\end{tabular}\bigskip\\
they are not needed.
If \twd{MV} (or \twd{CV}) is set then the external (vacuum) poloidal
field is added to the field of the plasma current.
In this case, the latter can also be zero.
All current densities in Tables 4.23 and 4.24 one are assumed to
be parallel to the magnetic field $\bf B$ 
\linebreak 
\indent
{\bf Table 4.24. Right hand sides of transport equations}\bigskip\\
\begin{tabular}{|c|l|c|l|}					\hline
\hspace*{-1pt}\bf Type\hspace*{-1pt} & \bf Variable 
			&\bf Units &\bf Description	\\	\hline
+ &\tw{CUBM} &${\rm MA}\,{\rm m}^{-2}$&  NB driven current\\
+ &\tw{CUECR}&${\rm MA}\,{\rm m}^{-2}$&  EC driven current\\
+ &\tw{CUFI} &${\rm MA}\,{\rm m}^{-2}$&  Fast ion current\\
+ &\tw{CUFW} &${\rm MA}\,{\rm m}^{-2}$&  FW driven current\\
+ &\tw{CUICR}&${\rm MA}\,{\rm m}^{-2}$&  IC driven current\\
+ &\tw{CULH} &${\rm MA}\,{\rm m}^{-2}$&  LH driven current\\	\hline
+ &\tw{PRAD / PRADX}& ${\rm MW}\,{\rm m}^{-3}$ & Radiation power\\
+ &\tw{PEECR}&${\rm MW}\,{\rm m}^{-3}$ & Electron heating due to ECR\\
+ &\tw{PEFW} &${\rm MW}\,{\rm m}^{-3}$ & Electron heating due to FW\\
+ &\tw{PEICR}&${\rm MW}\,{\rm m}^{-3}$ & Electron heating due to ICR\\
+ &\tw{PELH} &${\rm MW}\,{\rm m}^{-3}$ & Electron heating due to LHCD\\
+ &\tw{PIFW} &${\rm MW}\,{\rm m}^{-3}$ & Ion heating due to FW\\
+ &\tw{PIICR}&${\rm MW}\,{\rm m}^{-3}$ & Ion heating due to ICR\\
+ &\tw{PBEAM}&${\rm MW}\,{\rm m}^{-3}$ & Total beam power source\\ \hline
+ &\tw{SCUBM}& ${\rm kg}\;{\rm m/s}^2/{\rm m}^3$ 
				& Toroidal momentum source due to NB \\
+ &\tw{SNEBM}& $10^{19}{\rm\,m}^{-3}{\rm\,s}^{-1}$ 
				& Electron source due to NB\\
+ &\tw{SNNBM}& $10^{19}{\rm\,m}^{-3}{\rm\,s}^{-1}$ 
				& Warm neutral source due to NB\\
+ &\tw{SNIBMj}& $10^{19}{\rm\,m}^{-3}{\rm\,s}^{-1}$ 
	& Source of ions with the energy \tw{EBEAM/j, j=}1,2,3\\ \hline
\end{tabular}\bigskip\\
because they are supposed to be used in the parallel Ohm law \req{Ohm}. 
The only exception is \twd{CUTOR} which participates in the Joule
heating term \req{QJ}.

Some frequently used quantities are given in Table~4.25.
Density and temperature distributions of the `cold' neutral atoms are
calculated by the subroutine \twd{NEUT} as described 
\linebreak	
\indent{\bf Table 4.25. Auxiliary plasma parameters}\bigskip\\
\begin{tabular}{|c|l|c|l|}				\hline
\hspace*{-1pt}\bf Type\hspace*{-1pt} & \bf Variable 
&\bf Units &\bf Description	\\			\hline
$*$ &\tw{VP}     & m/s 
	    & $U_\pll/(d\psi/d\rho)=E_\pll/B_{\rm p},$ pinch velocity\\
$*$ &\tw{VPFP}  & m/s & $0.5\rho\dot B/B,$ flux surface velocity \\
+ &\tw{ER}    & V/m & Radial electric field\\
+ &\tw{VPOL~/~VPOLX}  & m/s & Poloidal rotation velocity\\
+ &\tw{VTOR~/~VTORX}  & m/s & Toroidal rotation velocity\\	\hline
+ &\tw{NN} & -- & Relative neutral density, $N(\rho)/N(\rho_B)$\\
+ &\tw{TN} & keV & Neutral temperature, $T_N(\rho)$	\\	\hline
+ &\tw{PBLON}& $10^{19}{\rm\,keV/m}^{3}$ 
			& Longitudinal pressure of  fast ions\\
+ &\tw{PBPER}& $10^{19}{\rm\,keV/m}^{3}$
			& Perpendicular pressure of  fast ions\\
+ &\tw{PELON}& $10^{19}{\rm\,keV/m}^{3}$ 
			& Longitudinal electron pressure\\
+ &\tw{PEPER}& $10^{19}{\rm\,keV/m}^{3}$ 
			& Perpendicular electron pressure\\	\hline
\end{tabular}\bigskip\\%\vskip 3ex
in Section~4.11.

Table~4.26 contains parameters characterizing the tokamak magnetic
configuration. 
The quantities $\tt EQFF$ and  $\tt EQPF$ give the right hand side of
the Grad-Shafranov equation \req{GSh}.
All others are obtained as a result of solving this equation.
Different average magnetic field characteristics are used in transport
coefficients while purely geometric quantities are included in the
transport set of equations as metric coefficients.
\\[1ex]
\indent
{\bf Table 4.26. Geometric and magnetic field characteristics}\bigskip\\
\begin{tabular}{|c|l|c|l|}					\hline
\hspace*{-1pt}\bf Type\hspace*{-1pt} & \bf Variable 
		&\bf Units &\bf Description	\\		\hline
+ &\tt EQFF~ / EQPF  & MA$\,\rm m^{-2}$ & $j_\zeta(r,z)= 
		R_0/r*\tw{EQFF} + r/R_0*\tw{EQPF}$ \\		\hline
+ &\tt BMAXT / BMINT & T & Maximum / minimum B field on a surface\\
+ &\tt B0DB2 / BDB02 & & $<B_0^2/B^2>$ \tw{~/~}	$<B^2/B_0^2>$\\
+ &\tt BDB0 & & $<B/B_0>$\\
+ &\tt FOFB  & & $<B_0^2/B^2(1.-\sqrt{1-B/B_{max}}(1+.5B/B_{max})>$ \\
								\hline
+ &\tt GRADRO& & $<|\nabla\rho|>,$ lateral area $S_a=$\tw{VR*GRADRO}\\
+ &\tt DRODA / DRODAX & &   d$\rho$/da\\
+ &\tt G11~~ / G11X & m$^2$ & $<(\nabla\rho)^2>V'$\\
+ &\tt G22~~ / G22X & m~ & $V'R_0/(4\pi^2 J)<(\nabla(\rho)/r)^2>$\\
+ &\tt G33~~ / G33X & & $<R_0^2/r^2>$\\
+ &\tt IPOL~ / IPOLX&& $J,$ Normalized poloidal current \req{dfJ}\\
								\hline
\end{tabular}\vskip 3ex

The next group of parameters (Table~4.27) describes plasma isotope
composition and impurity species. 
These quantities are used in subroutines 
{\tt ~NCLASS, NBI, STRAHL~} and in many other code modules
dealing with the transport properties, additional heating and so on.
\\
\indent
{\bf Table 4.27. Plasma isotope and impurity composition}\bigskip\\
\begin{tabular}{|c|lll|c|l|}
\hline
\hspace*{-1pt}\bf Type\hspace*{-1pt} &\multicolumn{3}{|c|}{\bf Variable} 
		&\bf Units &\bf Description		\\	\hline
+ &\tt NHYDR &\tt NDEUT &\tt NTRIT 	
	& $10^{19}$${\rm m}^{-3}$	& Hydrogen isotope density\\
+ &\tt NHE3  &\tt NALF  &  	
	&$10^{19}$${\rm m}^{-3}$ & Helium isotope density\\
+ &\tt NIZ1  &\tt NIZ2  &\tt NIZ3 
	& $10^{19}$${\rm m}^{-3}$ 	& Impurity species density\\
+ &\tt ZIM1  &\tt ZIM2  &\tt ZIM3 
	& $e_p$ 		& Impurity species charge \\ 
+ &\tt ZEF1  &\tt ZEF2  &\tt ZEF3
	& 			& Impurity contribution to $Z_{eff}$\\
+ &\tt DIMP1 &\tt DIMP2 &\tt DIMP3
	& ${\rm m}^2/{\rm s}$ 	& Diffusion coefficient \\
+ &\tt VIMP1 &\tt VIMP2 &\tt VIMP3
	& m/s 			& Impurity pinch velocity \\
+ &\tt PBOL1 &\tt PBOL2 &\tt PBOL3
	& ${\rm MW}\,{\rm m}^{-3}$	& Contribution to bolometric power \\
+ &\tt PSXR1 &\tt PSXR2 &\tt PSXR3
	& ${\rm MW}\,{\rm m}^{-3}$ 	& Contribution to soft X-ray radiation \\
\hline
\end{tabular}\bigskip

The group of variables used in the additional transport equations
\rf{Aux}, \rf{QFj} is given in Table~4.28.
They are analogous to the variables used for the main plasma parameters
such as \twd{NE} or \twd{TE}.
\medskip\\
\indent
{\bf Table 4.28. Arrays associated with dummy unknowns \twd{F1, F2, F3}}
\\[2ex]
\begin{tabular}{|c|lll|c|l|}					\hline
\hspace*{-1pt}\bf Type\hspace*{-1pt}&\multicolumn{3}{|c|}{\bf Variable} 
	&\bf Units 	&\bf Description	\\		\hline
+ &\tt F1  & \tw{F1O} & \tw{F1X} & [$f_1$] & 
	$f_1(t)$ \tw{~/~} $f_1(t-\tau)$ \tw{~/~} $f_1^{exp}(t)$ \\
+ &\tt F2  & \tw{F2O} & \tw{F2X} & [$f_2$] & 
	$f_2(t)$ \tw{~/~} $f_2(t-\tau)$ \tw{~/~} $f_2^{exp}(t)$ \\
+ &\tt F3  & \tw{F3O} & \tw{F3X} & [$f_3$] & 
	$f_3(t)$ \tw{~/~} $f_3(t-\tau)$ \tw{~/~} $f_3^{exp}(t)$ \\
+ &\tt DF1 & \tw{DF2} & \tw{DF3} & ${\rm m}^2/{\rm s}$ & 
			Diffusion for the variable $f_j$	\\
+ &\tt VF1 & \tw{VF2} & \tw{VF3} & m/s & 
			Pinch velocity for the variable $f_j$	\\
+ &\tt QF1 & \tw{QF2} & \tw{QF3} &[$f_j$]m$^3$/s & 
			Entire flux for the variable $f_j$	\\
+ &\tt SF1 & \tw{SF2} & \tw{SF3} &[$f_j$]/s & 
			Explicit rhs for the variable $f_j$	\\
+ &\tt SFF1&\tw{SFF2} & \tw{SFF3}& 1/s & 
			Implicit rhs for the variable $f_j$	\\
+ &\tt SF1TOT&\tw{SF2TOT}& \tw{SF3TOT} &[$f_j$]/s & 
			Total source for the variable $f_j$
%	\tw{SFjTOT=SFj+SFFj*Fj}	
\\	\hline
\end{tabular}\vskip 3ex %\bigskip\\

Finally, two groups of vector variables without predetermined
meanings are defined.
All can be used similarly to the scalar variables listed in
Table~4.17. 
Additionally, those of them with an \twd{`X'} at the end can be read in
from the data file (see Section~5.3).
\linebreak
\indent
{\bf Table 4.29. Dummy arrays}\bigskip\\
{\tt\begin{tabular}
{|p{15mm}p{15mm}p{15mm}p{15mm}p{15mm}p{15mm}p{15mm}p{15mm}|}	   \hline
CAR1	& CAR2	& CAR3	& CAR4	& CAR5	& CAR6	& CAR7	& CAR8	\\
CAR9	& CAR10	& CAR11	& CAR12	& CAR13	& CAR14	& CAR15	& CAR16	\\
CAR17	& CAR18	& CAR19	& CAR20	& CAR21	& CAR22	& CAR23	& CAR24	\\ 
CAR25	& CAR26	& CAR27	& CAR28	& CAR29	& CAR30	& CAR31	& CAR32	\\ \hline
CAR1X	& CAR2X	& CAR3X	& CAR4X	& CAR5X	& CAR6X	& CAR7X	& CAR8X	\\
CAR9X	& CAR10X& CAR11X& CAR12X& CAR13X& CAR14X& CAR15X& CAR16X\\ \hline
\end{tabular}}\vskip 3ex

\subsection{Astra expressions}

\subsubsection{Types of expressions}

Astra expressions provide a convenient notation for defining
physical quantities or discharge characteristics which can be expressed
by a formula. 
Each formula/expression can be referred to by a name and used in other
expressions as described in Section~5.2.
Typical examples are transport coefficients, power and particle sources
and sinks, cross-sections of atomic processes and so on.  
Once these expressions are programmed they are stored in an object
library or as Fortran text files.
The Astra compiler then ensures the appropriate treatment of all
modules and includes them into the resulting code according to their
meaning, implementation and application.
The whole concept is rather flexible: the expressions can be radially
and/or time dependent, they can include integrals and derivatives, depend
on other expressions and plasma parameters.
In many cases, the Astra expressions are self-documented because they
contain comments, examples of usage in the code, and references to 
papers from which they originate. 

On the one hand, the number of Astra expressions is continuously
increasing, on the other hand, some which have become obsolete are
retained for backward compatibility.
Therefore, instead of comprehensively describing all Astra expressions,
the goal of this section is to give a general representation of the Astra
expression library, to describe its structure, and to demonstrate how to 
find out more. 

Astra supports two different kinds of expressions:
formulae and functions. 
A beginning user of the code does not need to discriminate between
functions and formulae and can treat them as identical.
Moreover, the names of Astra expressions can coincide with the names of
the vector variables described in the previous section. 
The only essential difference for a user is that, unlike arrays,
expressions can be modified and new ones can be added. 
A user can delete any formula (including those in the Astra kernel) and
replace it with his own. 
A function provided with the code cannot be deleted, however, this is
not very restrictive because a user can always create his own function
with a new name.
Although no detailed explanation of the expression syntax is provided
here, a number of expressions in the code can be used as templates for
designing new ones. 

The actual difference formulae and functions is that the formulae are
included in a source code with the ${\tt INCLUDE}$ statement, while
functions are conventional Fortran functions. 
Although there is no clear distinction, functions tend to be used for
programming more complicated expressions than formulae.\footnote[1] 
{This sequence can be continued by the Astra subroutines which are
described in Section~4.9 and related to the processes described by
equations rather than by expressions.}
Nevertheless, many expressions are represented both by
functions and by formulae. 
Usage of formulae is preferable because it does not require rebuilding
libraries and usually results in a faster executable code than using
functions. 
As already mentioned, the proper usage is handled by the Astra compiler
so that a user need not go into these details.

In this description, only the overview and usage of Astra
expressions are presented. 
%Syntax and details will be discussed elsewhere.
Usually, a name of variable/expression has a mnemonic meaning
related to its use in the code. 
So the expressions supposed to be used as the matrix elements or the
right hand sides in Eqs.~\rf{Amain}--\rf{Afluxes} start with the same
letters as the corresponding terms of the equations. 
With a few exceptions the first letter of an expression name has the
following meaning
\bigskip \\ 
\begin{tabular}{l l}
\tw{B} &-- contributions to $\beta$'s,\medskip \\
\tw{C} &-- conductivity,\medskip \\
\tw{D} &-- diffusivity,\medskip \\
\tw{F} &-- predefined radial functions,\medskip \\
\tw{H} or \tw {X} &-- heat conductivity, (\tw{`H'} is used for electron, 
		\tw{`X'} for ion heat conductivity),\medskip \\
\tw{I} &-- surface integral (applied for current density),\medskip\\
\tw{L} or \tw{R} &-- plasma inductance or characteristic length,
							\medskip\\
\tw{N} &-- plasma collision frequencies,		\medskip\\
\tw{P} &-- heat source or sink,\medskip \\
\tw{Q} or \tw{W} &-- volume integral of a source or of energy
		content, respectively,\medskip \\
\tw{S} &-- particle source,\medskip \\
\tw{SV} &-- $<\sigma v>$ for collision processes,\medskip\\
\tw{V} &-- velocity.\medskip \\
\end{tabular}\\
More detailed information about every particular expression can be found
in its source file.
The source files for Astra formulae are placed in the directory
\twd{AWD/fml/}, functions in the directory \twd{AWD/fnc/}, arrays are
listed in the file \twd{AWD/for/status.inc}. 

\subsubsection{List of expressions}

In an Astra model (see Section~5.2) both upper- and lowercase letters can
be used. 
Below we write all Astra expressions in uppercase but emphasize that 
all file names for Astra expressions consist of lowercase letters only.
For instance, a source file for the expression \twd{QDT} appears as 
\twd{AWD/fml/qdt} (formula) and as \twd{AWD/fnc/qdt.f} (function).
Other expressions can be presented either by a formula or by a function
only. 
A user should also be aware that all variables associated with Astra
expressions are interpreted by the Astra code as Fortran real variables
independently of the first letter of the name.  
\smallskip\\
\noindent
{\bf Conductivity and CD efficiency}\\[1ex]
{\tt\begin{tabular}
{|p{13mm}p{13mm}p{13mm}p{13mm}p{13mm}p{13mm}p{13mm}p{13mm}p{13mm}|}\hline
CCSP & CCSPX & CNHH & CNHR & CCNEU & CCNEU & CHOTF & EFLHN & EFLHW\\ 
\hline
\end{tabular}}
\smallskip\\
\noindent
{\bf Bootstrap current density}\nopagebreak\\[1ex]
{\tt\begin{tabular}
{|p{13mm}p{13mm}p{13mm}p{13mm}p{13mm}p{13mm}p{13mm}p{13mm}p{13mm}|}
\hline
DCHA	& HCHA	& XCHA	& DCHH	& HCHH	& XCHH	& DCHR	& HCHR	& XCHR	\\
DCKIM	& HCKIM	& XCKIM	& DCKM1	& HCKM1	& XCKM1	& DHKIM	&	& \\
\hline
\end{tabular}}
\bigskip
\\
\noindent
{\bf Transport coefficients}\\[1ex]
{\tt\begin{tabular}
{|p{13mm}p{13mm}p{13mm}p{13mm}p{13mm}p{13mm}p{13mm}p{13mm}p{13mm}|}
\hline
CERL	& DBOHM	& DNDIF	& DNEXP	& ENHH0	& ENHH1	& ENHH2	& FOWC	& HAALC	\\
HABM	& HABMS	& HABOM	& HACTE	& HAED	& HAETI	& HAETS	& HAGB	& HAGBS	\\
HAGBS1	& HAITF	& HAITR	& HAMM	& HANAG	& HANAL	& HAPA	& HAPUE	& HAPYU	\\
HAQ1	& HAQ1C	& HARL	& HARLS	& HARNQ	& HARPL	& HASCL	& HATL	& HATLI	\\
HATLS	& HBJET	& HCHII	& HCHGP	& HEEFF	& HEGN 	& HETAI	& HETIS	& HEXP	\\
HGBEJ	& HGBIJ	& HIT89	& HNCHI	& HNGSB	& HNGSE	& HNGSI	& HNGSP	& HNPSI	\\
TECRL	& XCH86	& XCHR	& XEXP	& XIGN	& XIEFF	& XIRL	&	& \\ 
\hline
\end{tabular}}
\bigskip
\\
\noindent
{\bf Power sources and sinks}\\[1ex]
{\tt\begin{tabular}
{|p{13mm}p{13mm}p{13mm}p{13mm}p{13mm}p{13mm}p{13mm}p{13mm}p{13mm}|}   \hline
PAION	& PBDHE	& PBEIE	& PBICX	& PBRAD	& PDT	& PEDT	& PEDT1	& PEGN	\\
PEHCL	& PEI	& PEICL	& PEIGN	& PENEU	& PENLI	& PETSL	& PHICL	& PICX	\\
PIDT	& PIDT1	& PIGN	& PINEU	& PIONZ	& PIREC	& PITCX	& PITSL	& PJOUL	\\
POH	& PRCAR	& PRFER	& PRNEO	& PRNIT	& PROXI	& PROXY	& PRWOL	& PSYNC \\
\hline
\end{tabular}}
\bigskip
\\
\noindent
{\bf Particle sources and sinks}\\[1ex]
{\tt\begin{tabular}
{|p{13mm}p{13mm}p{13mm}p{13mm}|}				\hline
SNNEU	& SNNIE	& SNNII	& SNNR	\\			\hline
\end{tabular}}
\bigskip
\\
\noindent
{\bf Volume integrals of power and particle sources}\\[1ex]
{\tt\begin{tabular}
{|p{13mm}p{13mm}p{13mm}p{13mm}p{13mm}p{13mm}p{13mm}p{13mm}p{13mm}|} \hline
QBREM	& QDT	& QEDT	& QEGN	& QEICL	& QEIGN	& QENEU	& QETOT	& QEX\\
QIDT	& QIGN	& QINEU	& QITOT	& QIX	& QJOUL	& QNTOT	& QNX	& QOH\\
QRAD	& QRADX	& QSYNC	& QTOT	& QBTOT	& QEDWT	& QIDWT	& QNDNT	&\\ \hline
\end{tabular}}
\bigskip
\\
\noindent
{\bf Volume average quantities}\\[1ex]
{\tt\begin{tabular}
{|p{13mm}p{13mm}p{13mm}	p{13mm}	p{13mm}	p{13mm}	p{13mm}	p{13mm} p{13mm}|}
\hline
NEAV	& NEXAV	& TEAV	& TEXAV	& TIAV	& TIXAV	& TENDN	& TINDN &     \\
WALF	& WBPOL	& WE	& WEX	& WI	& WIX	& WTOT	& WTOTX	& ZNDN\\ 
\hline
\end{tabular}} 
\bigskip
\\
\noindent
{\bf Confinement times}\\[1ex]
{\tt\begin{tabular}
{|p{13mm}p{13mm}p{13mm}	p{13mm}	p{13mm}	p{13mm}	p{13mm}	p{13mm}|} \hline
TAUE	& TAUEE	& TAUEI	& TAUG	& TAUNA	& TAUP	& TAU89	& TITER\\ \hline
\end{tabular}}
\bigskip
\\
\noindent
{\bf Averaged cross-sections of atomic and nuclear processes}\\[1ex]
{\tt\begin{tabular}
{|p{13mm}p{13mm}p{13mm}p{13mm}p{13mm}p{13mm}p{13mm}p{13mm}|}	\hline
SVCX	& SVCXX	& SVD1	& SVD2	& SVDBH	& SVDHE	& SVDT	& SVIE	\\ 
SVIEP	& SVIEX	& SVII	& SVREC	& SVRECX& 	&	& \\	\hline
\end{tabular}}
\bigskip
\\
\noindent
{\bf Collisions}\\[1ex]
{\tt\begin{tabular}
{|p{13mm}p{13mm}p{13mm}p{13mm}p{13mm}p{13mm}p{13mm}|}		\hline
COULG	& NUE	& NUEE	& NUES	& NUI	& NUIS	& NUPP	\\	\hline
\end{tabular}}
\bigskip
\\
\noindent
{\bf Velocities}\nopagebreak\\[1ex]\nopagebreak
{\tt\begin{tabular}{|p{13mm}p{13mm}p{13mm}p{13mm}p{13mm}|}	\hline
CS	& VDIA	& VSI	& VTE	& VTI	\\			\hline
\end{tabular}}
\bigskip
\\
\noindent
{\bf Pressure and $\beta$'s}\\[1ex]
{\tt\begin{tabular}
{|p{13mm}p{13mm}p{13mm}	p{13mm}	p{13mm}	p{13mm}	p{13mm}|}	\hline
BETE	& BETI	& BETPL	& BETR	& BETAJ	& BETBM	& BETT	\\ 
ALMHD	& PRESE	& PRESI	& PREST	&	&	&	\\	\hline
\end{tabular}}
\bigskip
\\
\noindent
{\bf Integrals of current density}\\[1ex]
{\tt\begin{tabular}
{|p{13mm}p{13mm}p{13mm}	p{13mm}	p{13mm}	p{13mm}|}	\hline
IBM	& IBS	& ICD	& IOHM	& ITOT	& IX		\\
IECR	& IFI	& IFW	& IICR	& ILH 	&		\\	\hline
\end{tabular}}
\bigskip
\\
\noindent
{\bf Properties of the current profile}\\[1ex]
{\tt\begin{tabular}
{|p{13mm}p{13mm}p{13mm}p{13mm}p{13mm}p{13mm}p{13mm}p{13mm}p{13mm}|}
\hline
LICD	& LINT & QBM	& QBS	& QCD	& QNIND	& QMIN	& SHEAR	& XQMIN\\
\hline
\end{tabular}}
\bigskip
\\
\noindent
{\bf Characteristic lengths and derived quantities}\\[1ex]
{\tt\begin{tabular}
{|p{13mm}p{13mm}p{13mm}p{13mm}p{13mm}p{13mm}p{13mm}p{13mm}|}	\hline
LNE	& LNI	& LNZ1	& LTE	& LTI	& ETAN	& ETAE	& ETAI	\\
RLTCR	& RLTCZ	& RLTWN	& RLI	& RLS	&	&	&	\\ \hline
\end{tabular}}
\bigskip
\\
\noindent
{\bf Average Z in coronal approximation}\\[1ex]
{\tt\begin{tabular}
{|p{13mm}p{13mm}p{13mm}p{13mm}p{13mm}p{13mm}|}			\hline
ZICAR	& ZIFER	& ZINEO	& ZINIT	& ZIOXI	& ZIWOL \\		\hline
\end{tabular}}
\bigskip
\\
\noindent
{\bf Other plasma parameters}\\[1ex]
{\tt\begin{tabular}
{|p{13mm}p{13mm}p{13mm}	p{13mm}	p{13mm}	p{13mm}	p{13mm}	p{13mm} p{13mm}|}
\hline
EPAR	& EPL	& EDR	& D2TI	& SQZ	& DIDT	& FTE	& FTLLM	& TPF \\
ROTSH	& NECH  & NELA  & CCMHD	& CUOHM	& HMHD1	& HMHD2	& XMHD1	& XMHD2\\
\hline
\end{tabular}} 
\bigskip
\\
{\bf Radial functions} (See Section~4.8.3)\\[1ex]
{\tt\begin{tabular}
{|p{13mm}p{13mm}p{13mm}p{13mm}p{13mm}p{13mm}p{13mm}|}		\hline
FA	& FLIN	& FPA	& FPR	& FR	& FRS	&  FX	\\	\hline
\end{tabular}}

\subsection{Built-in functions}

Below we list some of the modules (Fortran subroutines and functions)
incorporated in the package. Headers of each module are equipped with
a brief description and examples of their usage. 
Fortran sources for most of the functions discussed in this section are
available in the file \twd{AWD/for/intern.f.}\footnote
{Changes in this file are not allowed for a non-authorized user.}
The general rules of calling Astra subroutines and functions are
described in Section~5.2.

The following notations for the arguments of Astra functions and
subroutines are adopted below:\\
\hspace*{1em}
$R$ is a vector variable, $R=R(\rho,t),$\\ \hspace*{1em}
$S$ is a scalar variable or an Astra vector at a given radius,
$S=S(t),$ \\ \hspace*{1em}
$G$ is a generic name for both, $R(\rho,t)$ and $S(t)$,\\ \hspace*{1em}
$t_0$ is the initial time of the simulation run,\\ \hspace*{1em}
$t\;$ is the current time, $t\geq t_0,$\\ \hspace*{1em}
$t_1$ is arbitrary time.\\
In the case where a built-in function allows a generic argument $G$ the
result has the same type as the argument. 
As a rule, only Astra variables (scalars and vectors) can be used as
arguments of Astra built-in functions.
To overcome this restriction one has to perform two steps as
in the following example:
\\ \indent
\tw{CAR1=WTOT;~~~~~~CAR2=TIMDER(CAR1)}
\\
Because \twd{WTOT} is an Astra formula, rather than an array, it cannot
be directly used as an argument of the built-in function \twd{TIMDER}
(time derivative).
Therefore, the first statement assigns the kinetic energy density 
$W_{\rm kin}(\rho)$ to the vector variable \twd{CAR1}, 
then the second computes the energy derivative as the vector~\tw{CAR2}:
$$
\tw{CAR1}=\tw{WTOT}=W_{\rm kin}
=0.0024\int_0^{\rho}\!\lb(n_eT_e+n_iT_i\rb)V'd\rho
\quad\rm[MJ]
,\qquad
\tw{CAR2}=\prti{W_{\rm kin}}{t}
\quad\rm[MW].
$$
In the next example 
\\ \indent
\tw{CV1=TIMDER(IPL);~~~~~~CV2=TIMINT(QNB)}
\\
both results are scalars, where ${\tw{CV1}=dI_{\rm pl}/dt}$
is the plasma current ramp rate and
$\tw{CV2}=\int_{t_0}^t\tw{QN}(\rho_{\rm B},t)dt$ is
the total particle flux through the outermost magnetic surface.

The rule restricting the argument type has two exceptions.
First, the functions \twd{CUT}, \twd{STEP} (see the end of 
Section~4.8.2) and Fortran intrinsic functions (Section~4.8.3)
allow arbitrary expressions as arguments.
Secondly, the values of Astra arrays at fixed radii (e.g. \twd{QNB} in
the example above) can also be arguments of the built-in functions.

\subsubsection{Time dependent functions}

\paragraph{Time derivative} (function {\tt TIMDER}).
$$
{\tt TIMDER}(G)=\lb.\lb(\prti{G}{t}\rb)\rb|_\rho
.$$

\paragraph{Time integral} (function {\tt TIMINT}).
$$
{\tt TIMINT}(G)=\int\limits_{t_0}^t G([\rho,]\,t)dt
.$$
Usage of functions \twd{TIMDER} and \twd{TIMINT} has an additional
limitation. 
They can only be used in an Astra model on the right hand side of an
assignment statement. 
This means that the constructions 
\\ \indent
\tw{CV1=TIMDER(VOLUME);~~~~~~dVdt\symbol{95}CV1}
\\
and
\\ \indent
\tw{CAR1=TIMDER(FP);~~~~~~~~~U\symbol{95}pl\symbol{92}CAR1}
\\
are correct while the result of \twd{dVdt\symbol{95}TIMDER(VOLUME)}
will be wrong.

\paragraph{Heaviside function of time} (function {\tt FJUMP}).
$$
{\tt FJUMP}(t_1)=H(t-t_1)
\lb\{\begin{array}{ll}
0
,\quad &
\mbox{if}\quad t<t_1\\
1
,&
\mbox{if}\quad t_1<t
.\\
\end{array}\rb.
$$

\paragraph{Linear ramp in time} (function {\tt FRAMP}).
$$
{\tt FRAMP}(t_1,t_2)=
\lb\{\begin{array}{ll}
0
,\quad &
\mbox{if}\quad t\leq t_1\\
(t-t_1)/(t_2-t_1)
,&
\mbox{if}\quad t_1< t< t_2\\
1
,&
\mbox{if}\quad t_2\leq t
.\\
\end{array}\rb.
$$

\paragraph{Storing of a variable value} at a given time 
(function {\tt FIXVAL}).
$$
{\tt FIXVAL}(G,t_1)=
\lb\{\begin{array}{ll}
G(t)
,\quad &
\mbox{if}\quad t<t_1\\
G(t_1)
,&
\mbox{if}\quad t_1<t
.\\
\end{array}\rb.
$$

\paragraph{Sliding time average.} 
Function of two parameters $\tws{FTAV}(G,\Delta t)$
provides smoothing of the variable $G(t)$ over the time interval 
$\Delta t$ so that  
$$
\lb.\lb.\vphantom{\frac{1}{2}}{\tt FTAV(G,}\Delta t)\rb|_t=
{\tt FTAV(G,}\Delta t)\rb|_{t-\tau}\exp\lb(-\frac{\tau}{\Delta t}\rb)
+\tws{G}(t)\lb[1-\exp\lb(-\frac{\tau}{\Delta t}\rb)\rb]
,$$
with $\tau$ being the time step. In particular,
$$
\lb\{\begin{array}{ll}
{\tt FTAV}(t)
\approx G(t)
, &
\mbox{if}\quad \tau \gg \Delta t\\
{\tt FTAV}(t)
\approx {\tt FTAV}(t-\tau)
,\quad &
\mbox{if}\quad  \tau \ll \Delta t
.\\
\end{array}\rb.
$$

\paragraph{Minimum/maximum value in time of a function $G$ 
({\tt FTMIN, FTMAX}).} 
$$
{\tt FTMIN}(G)=\min\limits_{t_0\leq\theta\leq t}\lb[G(\theta)\rb]
,\qquad
{\tt FTMAX}(G)=\max\limits_{t_0\leq\theta\leq t}\lb[G(\theta)\rb]
.$$
Here $t_0=\tt TSTART$ is the initial time and $t=\tt TIME$ is the
current time of the run. 

\paragraph{Examples:}
\begin{description}
\item[\quad\tw{FJUMP(1.)}] 
provides a time dependent function with the unity step at the time
$t=1$s. The actual width of the transition phase is no less than the
current time step $\tau=\tw{TAU}.$
\item[\quad\tw{FRAMP(1.,1.01)}] extends the transition phase up to 10 ms
(provided $\tau<10$ms, otherwise $\tau$).
\item[\quad\tw{MU\hspace*{3pt}-\hspace*{3pt}FIXVAL(MU,1.)}] 
returns the finite difference 
$\Delta\mu(\rho,t)=\mu(\rho,t)-\mu(\rho,1.),$ if $t>1$s, 
and 0~otherwise.
\item[\quad\tw{FTAV(UPL,.01)}]
returns a vector variable which is the loop voltage 
$U_{\rm pl}(\rho,t)$ averaged over the recent 10 ms (assuming that the time step 
$\tau\ll 10$ms). 
\item[\quad\tw{FTAV(UPLB,.01)}] is a scalar giving the average edge loop
voltage $U_{\rm pl}(\rho_{\rm B},t).$ 
\item[\quad\tw{FTMIN(MU)}] and \tw{~FTMAX(MU)~} give a range of
variation of $\tw{MU}=1/q$ during a run.
\item[!\hspace*{9pt}\tw{FTAV(HEXP,.01)}] this usage is forbidden
because \twd{HEXP} is not an Astra array. 
Two statements 
\twd{CAR1=HEXP;~CAR2=FTAV(CAR1,.01)} should be used instead.
%\item[\quad\tw{FTMAX(QOHB)}] is allowed and gives the
%maximum in time value of the total Ohmic power. However, a repeated call
%as \\ \hspace*{5em}
%\tw{Q(0)\symbol{95}FTMAX(QOHC)}
%;\qquad 
%\tw{Q(a)\symbol{95}FTMAX(QOHC)}\\
%results in an error.
\item[\quad\tw{TIMINT(UPL)}] is a vector variable which gives the
poloidal flux 
$\psi(\rho,t)=\psi(\rho,t_0)+\int_{t_0}^tU_{\rm pl}(\rho,t)dt$
up to an additive constant.
\item[\quad\tw{TIMDER(FPB)}] is a scalar variable and within numerical errors coincides with 
$U_{\rm pl}(\rho_{\rm B},t).$
\end{description}

\subsubsection{Radially dependent functions}

\paragraph{Linear radial functions} ({\tt FA, FX, FLIN, FR, FRS}).\\
All functions listed above are implemented as Astra formulae. 
The function {\tt~FR~}(array synonym{\tt~RHO}) gives the nodes of the
main transport radial grid~$\rho_j$ 
$$
\rho_j=
\lb\{\begin{array}{ll}
h(j-0.5)
, &
\mbox{if}\quad 1 \leq j < {\tt NA1}\\
\rho_B
,\quad &
\mbox{if}\quad j = {\tt NA1}
.\\
\end{array}\rb.
$$
Similarly, ${\tt FA}$ (array synonym {\tt~AMETR}) gives values of
the variable $a,$ i.e. $a_j=a(\rho_j)$ on this grid.
Most frequently used functions are defined on this grid.
Functions ${\tt FX}$ and ${\tt~FLIN}$ both return the normalized
radius $\rho_j/\rho_B$ which can also be obtained as
$\tw{RHO/ROC}.$ 

The intermediate grid in $\rho$ is given by the function
${\tt FRS}=jh.$
The functions listed in Table~4.5, all fluxes and some others are
defined on this latter grid.  

\paragraph{Parabolic profiles} ({\tt FPR, FPA}).
The Astra formulae
$${\tt FPR}=1-(\rho/\rho_B)^2
\qquad{\rm and}\qquad
{\tt FPA}=1-(a/a_B)^2
$$ 
describe parabolic profiles in variables $\rho$ and $a,$ respectively.

\paragraph{Gaussian distribution} 
(subroutine {\tt FGAUSS}, function {\tt GAUSS}). \\
The subroutine {\tt FGAUSS}$(\rho_0,\Delta_\rho,R)$ returns a radial
Gaussian profile in the variable $\rho$
$$
R_i=N\exp\lb[-\lb(\frac{\rho_i-\rho_0}{\Delta_\rho}\rb)^2\rb]
,$$
where the factor $N$ (of dimension ${\rm m}^{-3}$) is determined by the
normalization condition 
$$
\int_0^{\rho_B} R(\rho)dV=1
.$$
This means that the statement 
%\\ \hspace*{1em}
``\twd{FGAUSS(.5,.2,PEX):;}''
%\hspace*{1em}%\\
defines the vector variable \twd{PEX} as a Gaussian profile, centered in
$\rho_0=0.5$m with the width $\Delta\rho=0.2$m.
If the quantity \twd{PEX} is defined as a power, then the total power
deposited would be 1 MW.

The function {\tt GAUSS}$(\rho_0,\Delta_\rho)$ can be used similarly:
\marginpar{!!}$\tt PEX = GAUSS(0.5,0.2),$ 
however, the function \twd{GAUSS} can be called in a model only once.

\paragraph{Minimum/maximum value} of a vector variable $R(\rho)$ 
({\tt FRMIN, FRMAX}).
$$
{\tt FRMIN}(R)=\min\limits_{0\leq\rho\leq\rho_B}\lb[R(\rho)\rb]
,\qquad
{\tt FRMAX}(R)=\max\limits_{0\leq\rho\leq\rho_B}\lb[R(\rho)\rb]
.$$

\paragraph{Position of minimum/maximum} of a vector variable $R(\rho)$ 
({\tt RFMIN, RFMAX}).\\
The functions {\tt RFMIN, RFMAX} return 
$$
\begin{array}{ll}
{\tt RFMIN}(R)=\rho_{\min}
,\quad\mbox{where}\quad
R(\rho_{\min})=\min\limits_{0\leq\rho\leq\rho_B}\lb[R(\rho)\rb]
,\bigskip\\
{\tt RFMAX}(R)=\rho_{\max}
,\quad\mbox{where}\quad
R(\rho_{\max})=\max\limits_{0\leq\rho\leq\rho_B}\lb[R(\rho)\rb]
.\end{array}
$$

\paragraph{Root of the algebraic equation} $R(\rho)=R_0$ ({\tt RFVAL}).
$$
{\tt RFVAL}(R,R_0)=\rho_0
,\quad\mbox{where}\quad
R(\rho=\rho_0)=R_0
.$$
Maximum of all possible solutions is returned, $\rho_0=0$ if no root
found. 

\paragraph{Root of the algebraic equation} $R(a)=R_0$ ({\tt AFVAL}).
$$
{\tt AFVAL}(R,R_0)=a_0
,\quad\mbox{where}\quad
R(a=a_0)=R_0
.$$
Maximum of all solutions is returned, 0 if no root found.

\paragraph{Gradient} of a function $R(\rho)$ ({\tt GRAD}).
$$
{\tt GRAD}(R)=\prti{R(\rho)}{\rho}
.$$
One should beware of the approximate character of this operation.
The function \twd{GRAD} is intended for auxiliary purposes such as
graphical presentation or estimates only.  
For instance, it makes no distinction between the main and intermediate
radial grids. 
Therefore, the safety factor $q$ calculated as 
\twd{GP2*BTOR*RHO/GRAD(FP)} shows significant deviation from the
exact (in a finite difference sense) quantity 
$q=2\pi B_0\rho\,\prt\rho/\prt\psi=1/\mu=\tt1/MU$ at the
plasma centre. 
Even the improved evaluation $\tt GP2*BTOR*FRS/GRAD(FP)$ is not
everywhere exact because it does not take into account the centering of
different grid variables appropriately. 

\paragraph{Volume integral} of a function $R(\rho)$ ({\tt VINT}).
$$
{\tt VINT}(R)=\int\limits_0^{\rho} R(\rho)V'd\rho
.$$
In Astra models the function \twd{VINT} can be used in several different
ways: 
$$
\tw{
CAR1=VINT(NE);    \qquad
CV1=VINT(NEB);    \qquad
CF1=VINT(NE,0.5);
}
$$
The first statement in this example produces a vector, others
scalar output.
Note that the `0.5' as the second argument of \twd{VINT} (in the 3rd 
statement) refers to the minor radius $a$ measured in meters and 
introduced by \req{3mom}.
Therefore, this gives the following results:
$$
\tw{CAR1}=\!\int\limits_0^{\rho}\tw{NE}(\rho)V'd\rho
,\qquad
\tw{CV1}=\!\int\limits_0^{a_B}\tw{NE}(a)dV
=\!\int\limits_0^{\rho_B}\tw{NE}(\rho)V'd\rho
,\qquad
\tw{CF1}=\!\!\int\limits_0^{0.5\,\rm m}\!\tw{NE}(a)dV
.$$
For instance, the density averaged over a volume inside the current
radius $\rho$ can be submitted to the radial graphics output by the two
instructions: 
$$
\tw{CAR1=1;
\qquad
<ne>\symbol{92}VINT(NE)/VINT(CAR1);~~~
}
$$
In the last example, \twd{VINT(CAR1)} is equivalent to \twd{VOLUM}.
Moreover, the entire expression for $<n_e>$ can be obtained making use
of already defined Astra function \twd{NEAV}.

Exactly the same rules are applicable to the next Astra function
\twd{IINT} which also calculates an integral over the poloidal
cross-section rather then over the volume of a magnetic surface.

\paragraph{Surface integral} of a function $R(\rho)$ in a poloidal plane
({\tt IINT}). 
$$
{\tt IINT}(R)
=\frac{J}{2\pi R_0}\int\limits_0^{\rho} R(\rho)\frac{V'}{J^2}d\rho
.$$
This function is used for integration of the longitudinal
current density (see~\req{Ipl}). 
Similar to \twd{VINT} it allows expressions like \twd{IINT(CUECR),}
\twd{IINT(CUECRB)} and \twd{IINT(CUECR,0.2).}

As a general remark, we note that constructions involving \twd{VINT} (or
\twd{IINT}) are rarely needed because Astra expressions are
available for most physical quantities relevant to transport analysis.
We also remind the reader that \twd{VINT(1)} or \twd{VINT(NE*TE)} are
not allowed because the argument should be a radial array (Astra vector
variable). 

\paragraph{Heaviside function} of the radial variable $a$ ({\tt ASTEP}).
$$
{\tt ASTEP}(a_0)=H(a-a_0)
$$

\paragraph{Heaviside function} of the radial variable $\rho$ ({\tt
RSTEP}). 
$$
{\tt RSTEP}(\rho_0)=H(\rho-\rho_0)
$$

\paragraph{Heaviside function} of the radial variable $x=\rho/\rho_B$
({\tt XSTEP}). 
$$
{\tt XSTEP}(x_0)=H(\rho/\rho_B-x_0)
$$

\paragraph{Heaviside function} of any dimensionless argument $A$ ({\tt
STEP}).
$$
{\tt STEP}(A)=H(A)
=\lb\{
\begin{array}{l}
1,\qquad\mbox{if}\quad A\geq 0,\\
0,\qquad\mbox{if}\quad A< 0.
\end{array}\rb.
$$
An argument $A$ of the function \twd{STEP} can be any expression allowed
by the Astra rules. The same is valid also for the next function.

\paragraph{Cut-off function ({\tt CUT}).}
Sometimes, a function for the graphic output varies by many order of
magnitude. 
Because of a very limited word length (1 byte) of the data stored in 
an Astra post-view file (Section~5.5.4) essential information can be
lost and saved data will have very low resolution. 
To avoid this one can cut off too large values of no interest:
$$
{\tt CUT}(a_0,A)=
\lb\{\begin{array}{lll}
-a_0, 	&\quad\mbox{if}\quad & \;\;\;A\leq -a_0,\\
A,	&\quad\mbox{if}\quad & -a_0\leq A\leq a_0,\\
a_0,	&\quad\mbox{if}\quad & \;\;\;a_0\leq A.\\
\end{array}\rb.
$$

\begin{description}
\item[Examples:]
\item[\quad\tw{RSTEP(.5)}] can also be obtained as \twd{STEP(RHO-.5)},~ 
\twd{XSTEP(.2)}~as~\twd{STEP(RHO/ROC-.2)},\linebreak
\twd{ASTEP(.1)}~is equivalent to~\twd{STEP(AMETR-.1)}. 
Note that the numerical argument of \twd{XSTEP} (0.2 in this example) is
dimensionless, all others are measured in meters.
\item[\quad\tw{CUT(100.,NUES)}] can be useful for drawing a radial
profile of the collision frequency which can be very large near the
magnetic axis and at the plasma edge.
The same effect can be achieved in a more complicated way as
\twd{100-(100-NUES)*STEP(100-NUES).}
\item[\quad\tw{STEP(SIN(GP2*100*TIME)-0.5)}] shows another example of
allowed usage for this function which gives a time dependent square-wave
signal with a period 0.01 s (frequency ${f=100}$Hz) and a duty factor 
$(\pi/2-\arcsin(0.5))/\pi=33\%$.
\item[\quad\tw{FRMAX(MU)}] returns the minimum value of the safety factor
$q(\rho)=1/\mu(\rho)=\tt 1/MU,$
\item[\quad\tw{RFMAX(MU)}] returns the $\rho\,-$position of this minimum,
\item[\quad\tw{RFVAL(MU,0.5)}] returns the position of the$q=2$resonance
surface in the coordinate$\rho,$
\item[\quad\tw{AFVAL(MU,0.5)}] returns the position of the same surface
with respect to the coordinate$a.$
The last result can also be obtained as 
%\twd{RADIAL(AMETRC,RFVAL(MU,0.5))} or shorter as
\twd{AMETR(RFVAL(MU,0.5)).}
\end{description}

\subsubsection{Functions for mapping one flux label to another}

Several Astra functions are available for recalculating one `radial' 
coordinate to another: {\tt RFA}, {\tt XFA}, {\tt AFR}, {\tt AFX}, 
{\tt RFAN}, {\tt XFAN} . 
The function \twd{RFA(0.5)} returns a value of the coordinate $\rho$
(in meters) at the magnetic surface labeled by the mid-plane radius
$a=0.5$m (see \req{3mom}).
The function \twd{XFA({\it a})} is equivalent to 
${\tt RFA}(a)/{\tt ROC}$ and returns the normalized value of the flux
label $x(a)=\rho(a)/\rho_B.$ 
The functions {\tt AFR} and {\tt AFX} are inverses of {\tt RFA} and
{\tt XFA}, respectively.
Finally, {\tt RFAN(0.5)} gives $\rho$ as a function of the dimensionless
argument $a_N=a/a_B=0.5,$ similarly, 
${\tt XFAN}(a_N)={\tt RFAN}(a_N)/\rho_B.$

Vector variable at the given radius can be computed using one of the two
functions {\tt ATX} and {\tt ATR}.
The difference between both is that the first one refers to the
normalized radius $\rho/\rho_B$ and the second to the dimensional radius
$\rho.$ 
For instance, \twd{ATX(TI,0.5)} returns a value of the ion temperature
at $\rho=0.5\rho_B$ while \twd{ATR(TI,0.5)} at $\rho=$0.5 m. 
Other options are discussed in Section~5.2.3.

\subsubsection{Fortran intrinsic functions}

In addition to special functions described above one can use also the
standard set of Fortran functions:
{\tt SIN,} {\tt COS,} {\tt TAN,} {\tt ASIN,} {\tt ATAN,} {\tt EXP,} 
{\tt ABS,} {\tt MIN,} {\tt MAX,} {\tt SQRT,} {\tt ALOG,} {\tt ALOG10,}  
{\tt ANINT,} {\tt REAL} and {\tt SIGN.}
The type of the argument can be anything.
In turn, it defines a type of the result so that {\tt~SIN(CV1)} is a
scalar, while {\tt~SQRT(NE*TE)~} and {\tt~ALOG(LTE)} are vectors. 

\subsection{Plug-in subroutines}

All subroutines described in this section can be found under the name 
\twd{AWD/sbr/xxx.f} where \twd{`xxx'} should be replaced by a name
referred below in the lowercase. 
As a matter of fact, the directory \twd{AWD/sbr/} includes many more
subroutines which can be useful in different applications.
We describe here only few of those.

%\subsubsection{Fourier transform in time ({\tt FOURV, FOURTR}).}
%The subroutines perform Fourier transform for Astra scalar variable and
%array, respectively.

\subsubsection{Ionization particle flux ({\tt GNEX, GNXSRC}).}

Both subroutines have no parameters and communicate with the Astra core
through common blocks.
The subroutine ${\tt GNEX}$ calculates the particle flux \twd{GNX} caused
by ionization of neutral atoms in the plasma volume and also due to
evolution of the electron density 
$$
{\tt GNX}
=\frac{1}{S_a}\int\limits_0^\rho\lb(V'S_e-\prti{V'n_e}{t}\rb)d\rho
,$$
where the toroidal surface area $S_a=\tt SLAT$ is defined by \req{Sa}.
It is assumed that the density of wall neutrals \twd{NN} and a source
due to the beam neutrals \twd{SNEBM} are calculated 
independently (for instance, by calling subroutines \twd{NEUT} and
\twd{NBI}). 
{\tt GNXSRC~} gives the same quantity but without allowance for the
neutrals from NBI. 

The quantity \twd{GNX} is {\em always} included in \req{Amain}. 
It can be useful when the particle diffusion equation is switched off
and the density evolution is taken from an experiment and is a
prescribed function of time and radius. 
Then the `experimental' particle flux \twd{QN=SLAT*GNX} will be included
in the power balance in accordance with \req{Amain},\rf{Afluxes} unless 
$\gamma_e=\gamma_i=0.$ 
On the other hand, the user should be aware that calling any of the
subroutines {\tt GNEX} or {\tt GNXSRC} together with solving the
diffusion equation will cause a side effect and include the same
particle flux twice. 

\subsubsection{Upper and lower function boundaries (\twd{MINMAX}).}
A call to the subroutine ${\tt MINMAX}(R,R_{min},R_{max})$ returns two
quantities 
$$
R_{min}(\rho)
=\mathop{\rm min}\limits_{t_1\leq\theta\leq t_2}R(\rho,\theta)
\qquad\mbox{and}\qquad
R_{max}(\rho)
=\mathop{\rm max}\limits_{t_1\leq\theta\leq t_2}R(\rho,\theta)
.$$
Although the same result can be obtained by calling the two functions
${\tt FTMIN}$ and  ${\tt FTMAX}$ using \twd{MINMAX} is more flexible
because it finds maxima and minima in a prescribed time interval
$t_1<t<t_2.$ 
So the statement
\\ \indent
\tw{MINMAX(CU,CAR1,CAR2)::0.5:0.6}
\\
gives range of variation of the current density for 100 ms inside the
time interval 0.5 s$<t<$0.6 s.

\subsubsection{Sawtooth oscillations (internal disruption) 
{\rm (subroutine \tt MIXINT)}.}

This subroutine models profile evolution of 
$T_e,T_i,n_e,n_i,j_\pll,\mu,\psi$ 
in the process of magnetic field line reconnection according to
Kadomtsev's theory of sawtooth oscillations \cite{BB}.
The current version of the code allows for the internal reconnection at
the resonance $q=m/n=1$ only.
In a process of the full reconnection final distributions of all plasma
parameters are uniquely determined by the conservation laws and by the
initial distributions. 
The only physical characteristic missing in  this theory is the trigger
of the internal disruption.
Therefore, the subroutine \twd{MIXINT} requires one physical input
parameter which can be either the sawtooth period or the reconnection
radius.
Additionally, it is possible to use a disruption condition
based on the double-tearing-mode reconnection model \cite{PP}.
In this case no additional parameters are required.
A statement calling the subroutine reads
\\ \indent
\tw{MIXINT(OPTION,RECON):}
\\
Both input parameters must be given as real Fortran variables or
constants and have the following significance:
\begin{description}
\item
$\tw{OPTION}=0:$ if a reconnection takes place at $t=t_1$ then the
next one occurs not earlier than at $t=t_1+\tt RECON$ and as soon as
the resonance surface $q=1$ is found,
\item
$\tw{OPTION}=1:$ a reconnection occurs as soon as the radius $\rho_1$
of the resonance surface $q=1$ satisfies the condition
$\rho_1/\rho_{\rm B}>\tw{RECON},$ 
\item
$\tw{OPTION}=2:$ a reconnection occurs according to the model of
\cite{PP}. 
The second parameter $\tt RECON$ is ignored.
\end{description}
The header of the subroutine \twd{MIXINT} describes other options which
provide extended information about every sawtooth event.

Reconnection mixes up all plasma parameters, however, in the code,
only those quantities are redistributed which evolution is described by
a transport equation.
For instance, if the equation for \twd{NE} is not solved then the
profiles of \twd{NE} and \twd{NI} do not change in the course of the
reconnection described by \twd{MIXINT}.
Similarly \tw{~FP, MU, CU~} do not change if the current density
distribution is prescribed.

\subsubsection{Gas puff neutrals {\rm (subroutine \tt NEUT)}.} 

The subroutine \twd{NEUT} solves an equation for the distribution of
`cold' (to distinguish from `fast' neutral particles from NBI) neutral
atoms in the plasma volume.  
The solution method first replaces the differential kinetic equation
\rf{fN} with an integral equation \cite{DK} then the integral equation
is solved iteratively. 
The parameter input/output for the subroutine can be schematically
depicted as 
$$
\begin{array}{lcr}
\fboxrule=1pt			\fboxsep=5mm
\fbox{\Large\bf Transport}
&
\begin{array}{c}
a_B,\;A_i,\;Z_i,\;n_e,\;T_{e,i};\;N_{1,2},\;E_{1,2},\;N_{iter}\\
\vector( 1,0){190}\\[-10pt]
\vector(-1,0){190}\\
N(\rho),\;T_N(\rho),\;\alpha_{\rm pl}
\end{array}
&
\fboxrule=1pt			\fboxsep=5mm
\fbox{\Large\tt NEUT}
\end{array}
$$
Most of the exchange parameters were already discussed.
Those which are directly related to the neutral atom distribution are 
listed in the following table.
\\[2ex]
%\noindent
\begin{tabular}{| l | l | l | l |}				\hline
\bf Parameter\hspace*{-4pt}&\bf Code name&\bf Units &\bf Description\\
								\hline
$N_1$	& {\tt NNCL}	&10$^{19}\rm m^{-3}$\hspace*{-7pt}
	& Density of incoming neutrals (1st component)\\
$N_2$	& {\tt NNWM}	&10$^{19}\rm m^{-3}$\hspace*{-7pt}
	& Density of incoming neutrals (2nd component)\\ 
$E_1$	& {\tt ENCL}	&keV	
	& Energy of incoming neutrals (1st component)\\
$E_2$	& {\tt ENWM}	&keV	
	& Energy of incoming neutrals (2nd component)\\
$N_{iter}$&\tt NNCX	& & Number of iterations in {\tt NEUT}\\ \hline
$N(\rho)$&\tt NN*(NNCL+NNWM)\hspace*{-4pt}&10$^{19}\rm m^{-3}$
		\hspace*{-7pt}&Density of neutral particles\\ 
$T_N(\rho)$&{\tt TN}	&keV& Temperature of neutral particles\\ 
$\alpha_{\rm pl}$&{\tt ALBPL}&	& Plasma albedo\\ \hline
\end{tabular}
\bigskip \\
The subroutine allows for two mono-energetic components of incoming
neutral particles which are characterized with their densities 
$N_1,N_2$ and energies $E_1,E_2$ (\twd{NNCL,NNWM} and \twd{ENCL,ENWM}),
respectively. 
It is assumed that the neutral mass coincides with a mass of the working
gas given by the variable \twd{AMJ}.
A control parameter \twd{NNCX} (default value 20) describes a number of
iterations for solving the integral equation.

Output of the subroutine \twd{NEUT} provides the normalized neutral
density \twd{NN} and the temperature \twd{TN}
distributions of the neutral particles in plasma and, additionally, the
plasma albedo \twd{ALBPL.}  
Optionally, it is also possible to calculate a distribution function of
the outgoing neutral particles. 
The quantity \twd{NN} is normalized in such a way
that, at the plasma edge, the density of incoming neutrals is
$N_1+N_2=\twd{NNCL+NNWM}$ and the density of outgoing neutrals is
$N_{\rm B}-N_1-N_2=\tt(NNCL+NNWM)*(NNB-1),$ where
$N_{\rm B}=N(\rho_{\rm B})$ and $\tw{NNB}=\tw{NN}(\rho_{\rm B}).$ 
Consequently, the average energy of incoming neutrals is
$$
E_{\rm in}=\frac{N_1E_1+N_2E_2}{N_1+N_2}
=\frac{\tt NNCL*ENCL+NNWM*ENWM}{\tt NNCL+NNWM}
$$
and outgoing
$$
E_{\rm out}=\frac{N_{\rm B}T_{\rm NB}-N_1E_1-N_2E_2}{N_{\rm B}-N_1-N_2}
=\frac{\tt NNB*TNB-NNCL*ENCL-NNWM*ENWM}{\tt(NNCL+NNWM)*(NNB-1)}
.$$
Finally, the inward flux is given by
$$
\Gamma_{\rm in}
=4.44\x10^{24}\ds\frac{\lb(N_1\sqrt{E_1}+N_2\sqrt{E_2}\rb)}{\sqrt{A_i}}
,\qquad\rm m^{-2}s^{-1}
$$
and the outward flux by
$$
\Gamma_{\rm out}=\alpha_{\rm pl}\Gamma_{\rm in}
.$$
\subsubsection{Density adjustment ({\tt SETNAV}).}

This subroutine adjusts the density of incoming cold neutrals in order to
provide a required volume average electron density. 
The following algorithm is adopted.
At each time step the particle confinement time $\tau_n$ is evaluated 
$\tau_n=\,<n_e>\!\!\lb/\lb(S-{d<n_e>}/{dt}\rb)\rb..$
Then the density of the cold neutrals (parameter \twd{NNCL}) is scaled
in order to provide $S_{\rm req}=<n_{e,\rm req}>/\tau_n.$ 
The resulting density evolution is approximately described by the 0D
equation 
$$
\frac{d<n_e>}{dt}
=\frac{<n_{e,\rm req}>-<n_e>}{\tau_n}+S-S_{\rm req}
,$$
where $S_{\rm req},\,\tau_n$ and $<n_{e,\rm req}>$ are taken from the
previous time step.
If initially the actual density $<n_e>$ is far from the requested value 
$<n_{e,\rm req}>$ the described procedure can give an unreasonably large
neutral influx which subsequently crashes the code. 
In order to avoid this, an additional control parameter $\tau_N$ is
introduced which restricts the rate of variation 
$\tw{NNCL}=N_1$ by the condition $N_1/(N_1/dt)<\tau_N.$
This parameter should be selected as $\tau_N\geq\tau_n.$
As another option, $\tau_N$ can be set negative. 
Then an automatic adjustment will be enabled which works reasonably if
the difference 
$|\!<\!n_{e,\rm req}\!>-<\!n_e\!>\!|$ is not too large.

\noindent{\bf Example:}

\hangindent=8mm\hangafter=0\noindent
The part of the model related to \twd{SETNAV} reads:
\\ \hspace*{15mm}
\tw{CF1=5+5*FRAMP(.5,1);~~~~~~!~Prescribed volume average density}
\\ \hspace*{15mm}
\tw{SETNAV(CF1,.01):;~~~~~~~~~!$\tau_N=$0.01 s}\\
%\hangindent=8mm\hangafter=0\noindent
This will result in a ramp of the electron density as prescribed by the
first statement. 

\subsubsection{Smoothing a function ({\rm subroutine \tt SMEARR}).}

The subroutine ${\tt SMEARR}(\alpha,R_{\rm in},R_{\rm out})$ finds a
vector function $R_{\rm out}(\rho)$ as an extremum of the functional 
$$
\mathop{\rm min}\int\limits_0^{\rho_B}\left[
\alpha\left(\prti{R_{\rm out}}{\rho}\right)^2
+\left(R_{\rm in}(\rho)-R_{\rm out}(\rho)\right)^2
\right]d\rho
$$
with the boundary conditions
$$
\left.\prti{R_{\rm out}}{\rho}\right|_{\rho=0}=0
,\qquad
R_{\rm out}(\rho_B)=R_{\rm in}(\rho_B)
$$
where $R_{\rm in}(\rho_j)$ is a given grid function, $\alpha$ is a
control parameter and $R_{\rm out}$ is the output function.
Reasonable values of the control parameter $\alpha$ should be around
0.001. 
The value $\alpha=0.01$ results in a very strong smoothing of the
result. 

\noindent{\bf Example:}
\\ \hspace*{2cm}
\tw{CAR1=HEEFF;~~~~~~SMEARR(0.002,CAR1,CAR2):;}
\\ \hspace*{0.5cm}
This gives a radially smoothed effective electron heat conductivity
\\ \hspace*{2cm}
$
\chi_{e,eff}=\int\lb(P_e-\prt W_e/\prt t\rb)dV/V'<(\na\rho)^2>n_e
$.

\subsubsection{Store array / variable value at a given time 
({\tt STARR / STVAR}).}

The subroutine ${\tt STARR}(R,t_1,R_1)$ (${\tt STVAR}(S,t_1,S_1)$)
stores a vector (simple) variable at a given time $t_1:$
$$
R_1(\rho)=R(\rho,t=t_1)
,\qquad\qquad
S_1=S(t=t_1)
.$$
Note that in recent versions of the code both subroutines can be
replaced by the function ${\tt FIXVAL}.$

\subsubsection{User defined time step control {\rm (subroutine \tt
TSCTRL)}.} 

The standard procedure of automatic time step control implemented in the
Astra code restricts the one-time-step variation in the main unknowns
\twd{TE, TI, NE, Fj}.
In some cases, for instance for stiff transport models, more
strict control is required.
 
Such a control can be achieved by calling the subroutine 
${\tt TSCTRL}(R_1,R_2,R_3,\Delta_{\rm max}).$ 
First, a variation in every vector input quantity $R_i,$ where $i=1,2,3,$ 
$$
\Delta_{Ri}=\mathop{\rm max}\limits_\rho
\left|\frac{R_i(\rho,t)-R_i(\rho,t-\tau)}{R_i(\rho,t)}\right|
$$
is evaluated.
Then the time step $\tau$ is multiplied by 
$\max\lb(\Delta_{\rm max}/\Delta_{Ri},\tau_{\rm inc}\rb).$
The quantity $\tau_{\rm inc}=\tt TAUINC$ is described in Table~4.15
and $\Delta_{\rm max}$ is the real input parameter which can vary
between 0.1 and 10. depending on the controlled quantities $R_i.$
In all cases, the time step $\tau$ is bounded by 
${\tt TAUMIN}<\tau<\tt TAUMAX.$ 

A limitation is that the subroutine ${\tt TSCTRL}$ can be called from a
model only once. 
To monitor several array parameters it allows up to three different
vector arguments.

\paragraph{Examples:}
\begin{description}
\item[$\,$1)] In a stiff transport model small variation in gradient
can cause huge variation in flux.
Therefore, plasma parameters like density or temperatures cannot provide
reliable time step control.
More appropriate regulation can be provided by \\
\hspace*{2cm}\twd{TSCTRL(QI,QE,QN,1.):;}
\item[$\,$2)] Another way to achieve a similar result is\\
\hspace*{2cm}\twd{TSCTRL(XI,XI,XI,CF1):;}\\
Here variation of a single quantity \twd{XI} is checked. 
The strength of the control can be adjusted interactively by changing
\twd{CF1}. 
\end{description}

\subsection{Interfaces to additional heating and CD modules}

\subsubsection{General features.}

The Astra system provides additional subroutines developed by different
authors for the main methods of additional plasma heating in tokamaks. 
Each of the subroutines is, in fact, a large package and we do
not pretend to describe them in this report. 
We discuss here the main features of the interfaces only, referring to
the original descriptions for more detail. 

\paragraph{Input.} 
The following set of input data is usually required 
for additional heating and current drive modules: 
\begin{itemize}
\item[]{\bf Configuration:}
$R_0${\tt (RTOR)}, 
$a_B${\tt (AB {\rm or} ABC)}, 
$a(\rho)${\tt (AMETR)}, 
$\Delta(\rho)${\tt (SHIF)}, \\
$\hphantom{{\bf Configuration:~}}\lambda(\rho)${\tt (ELON)}, 
$\delta(\rho)${\tt (TRIA)}, metric coefficients.
\item[]{\bf Plasma parameters:}
$A_j${\tt (AMJ {\rm or} AMAIN)}, 
$Z_{eff}(\rho)${\tt (ZEF)}, 
$Z_j${\tt (ZMJ {\rm or} ZMAIN,} \\
$\hphantom{{\bf Plasma parameters:}}$
{\tt ZIM1, ZIM2, ZIM3),} 
$n_{e,j}(\rho)${\tt (NE, NI, NHYDR, NDEUT,...} \\
$\hphantom{{\bf Plasma parameters:}}$
{\tt NIZ1, NIZ2, NIZ3)}, 
$T_{e,j}(\rho)${\tt (TE, TI).~~} Here index $j$\\
$\hphantom{{\bf Plasma parameters:}}$
extends over all ion species.
\item[]{\bf Current} and {\bf magnetic field:}
$I_{\rm pl}${\tt (IPL)}, 
$B_0${\tt (BTOR)}, 
$\mu(\rho)${\tt (MU)}, 
$j_\pll(\rho)${\tt (CU)}.
\item[]{\bf Launcher} or {\bf neutral beam parameters.}
\item[]{\bf Control parameters for the numerical algorithm.}
\end{itemize}

\paragraph{Output.} All modules return power deposition profiles,
$P_{e,i}^H(\rho),$ for electron and ion plasma components and, 
when applicable, driven current profile, $j_{CD}.$ 
Many other characteristics as injected momentum, distribution functions,
antenna coupling or shine through and so on are usually also available.

The set of input parameters can be split into two groups.
The first one includes those parameters, which are 
time-dependent or can change from shot to shot. 
The second group describes construction characteristics of a particular
tokamak or its heating system.
In most cases these parameters do not often change and only require be
set once. 
It is natural to exclude these parameters from the Astra core and
localize them inside the modules for additional heating.
The input parameters of the first group and the main output parameters
are treated as standard Astra control parameters which means that they
are accessible from the Astra code, can be defined, changed, plotted
using the standard Astra service as described in Sections~5.2 and 5.3.  

Different heating and current drive packages briefly described in
Section~3.14.1 can be obtained by special request.
The only exception is the NBI package which is always included in the
Astra code. 
For this reason the interface to the NB heating package is described in
more detail in the next Section.

\subsubsection{NB heating and current drive {\rm (subroutine \tt NBI \rm
by A.~R.~Polevoi)}} 

A neutral beam installation is represented in the NBI package by a set
of ``pencil beams''.  
Each pencil beam is characterized by its geometry and the power
distribution across the beam axis which in turn is described by a
number of thin parallel beams.
Presently, up to 16 pencil beams (or beam sources with the different
geometry) is allowed.

The NBI package considers (i) attenuation of the neutral beam during its
passage through a plasma due to ionization and charge-exchange, 
(ii) capture and losses (including ripple losses) of the new-born ions 
by analysis of the ion drift trajectories, 
(iii) thermalization of the supra-thermal ions and their
contribution to plasma heating, current drive and toroidal rotation. 
The latter problem requires solving the Fokker-Planck equation. 
In the fast standard version, an analytic solution of the steady-state
equation is used. A time dependent version is also available.

Interface between the Astra core and the neutral beam package is
provided by the subroutine \twd{AWD/sbr/nbi.f}.
As all other Astra plug-in modules this subroutine can be modified by
the user and, therefore, the description below should be viewed as a
framework providing rather guidelines than a comprehensive manual.

The usual set of general discharge and device parameters (see
Section~4.10.1) are transferred to the subroutine \twd{NBI} through
common blocks and is not discussed here. 
Additionally, a number of specific parameters describe the properties of
the NBI solver and the geometric characteristics of every pencil beam.
First of all, 8 control parameters 
\\
{\bf Table 4.30. Control parameters for the subroutine \twd{~NBI}}\\[2ex]
\begin{tabular}{|l|l|}					\hline
\bf Name	&\bf Description\\ \hline
\tt CNB1 	& A number of NBI sources (a number of pencil beams)\\
\tt CNB2 	& Explicit (=1) or implicit (=0) account for the charge
		  exchange losses\\ 
\tt CNB3 	& Space grid for NBI routine \twd{N1=(NA1-1)/CNB3+1}\\
\tt CNB4 	& Steady state (=1) or time dependent (=2) Fokker-Planck
		  solver \\
\tt CNBI1 	& Fast ion charge exchange losses due to cold neutrals\\
\tt CNBI2 	& Fast ion charge exchange losses due to NBI neutrals\\
\tt CNBI3 	& FP solver time step: $\tau_{NBI}=\tt CNBI2*TAU$\\
\tt CNBI4 	& FP solver velocity grid size: \twd{IV1=80/CNBI4+1}\\
\hline
\end{tabular}
\bigskip\\
are used to set different features of the NBI solver.
The control parameters are described in the Table~4.30 as they are used
in the subroutine \twd{AWD/sbr/nbi.f} although, in general, every
C-parameter (see Table~4.17) can be used to perform every control
function. 

The parameter \twd{CNB1} gives a number of pencil beams which have to be
taken into account. 
It can also take negative values which invoke an interactive mode for
modifying the beam geometry and will be discussed later.

The parameter \twd{CNB2} defines a way of describing the ion thermal
losses due to the charge exchange between the beam neutrals and plasma
ions.
These losses can be computed with the formula
\twd{PBICX=-0.0024*SNNBM*TI} where the source of thermal neutrals
\twd{SNNBM} is calculated by the NBI package. 
The corresponding losses can either be included into the ion
heating source \twd{PIBM} or calculated elsewhere.
If the parameter \twd{CNB2} is positive then the quantity \twd{PBICX}
is subtracted from the ion heating source \twd{PIBM}.
Otherwise, the losses should be taken into account in a model as shown
below 
\bigskip\\
%\begin{equation}\label{3models}
\begin{tabular}{p{5cm}|p{5cm}|p{4cm}}
\tt\quad Model 1 \rm (recommended)
		&\qq\qq\tt Model 2 &\qq\tt\quad Model 3 \\
\tt CNB2=0  	&\qq\tt CNB2=0  &\qq\tt CNB2=1  \\
\tt PI=PIBM+\dots &\qq\tt PI=PIBM-PBICX+\dots 
				&\qq\tt PI=PIBM+\dots \\
\tt PIT=-0.0024*SNNBM+\dots	& 		&
\end{tabular}\bigskip\\
%\end{equation}
Models 2 and 3 in this example are fully equivalent, however, the ion
heating term \twd{PIBM} is different in these two cases.
The model 1 enables implicit treatment of the charge exchange losses
which can result in more stable numerical scheme when these losses are
high enough.

The parameter \twd{CNB3} allows to run NBI routine on the reduced radial
grid in order to save computing time. 
This is reasonable when the fast ion drift trajectories are wide and
spread over many Astra radial grid steps.
Finally, the parameter \twd{CNB4} defines a type of the Fokker-Planck
solver to be used.
The C-parameters \twd{CNBI1$\div$CNBI4} are effective if \twd{CNB4=2}
only. 
Otherwise they can be anything.

As mentioned every beam source is represented by a pencil beam.
A sort of neutral atoms in the beam, their energy composition and power
distribution across the beam, geometry and some other characteristics
(altogether about 20 parameters) should be set separately for every
pencil beam. 
These parameter are stored in an NBI configuration file which can be
either attributed to every input data file (Section~5.1) or to every
experimental setup. 
%\footnote
%{For instance, four different NBI configuration files are
%used presently for the Asdex Upgrade tokamak:
%and two for the JET tokamak: \twd{AWD/exp/jetupshifted.nbi} and
%\twd{AWD/exp/jetstandard.nbi}.
%\twd{AWD/exp/aug99.nbi}, \twd{AWD/exp/aug00.nbi},} 
It is done in the following way. 
Assume that the Astra code is started with the data input file
\twd{aug10000}. 
When the \twd{NBI} routine is called the code checks whether the file
\twd{AWD/exp/aug10000.nbi} exists. 
If yes, then this file, otherwise, the default configuration file
\twd{AWD/exp/aug.nbi} will be used.
The latter can also serve as a template for creating new configuration
files. 

The NBI configuration file can start with arbitrary number of comment
lines. 
A comment line is a line with the sign \twd{'!'} in the first column.
The comment lines are allowed in the beginning of the
configuration file only.
A meaningful part of the file consists of \twd{CNB1} groups, one group
for every pencil beam. Each group consists of 21 fields. 
Each field occupies 12 positions and can be filled either with a name of
C-parameter (Table~4.17) or Z-parameter (Tables~4.18 and 4.19) or with a
number in the free format as shown in the example below:
\begin{equation}\label{NBI_conf}\begin{tabular}{l}\tt\hspace*{2ex}
1						\\ \tt%\hspace*{1ex}
~~ZRD1~~~~~~0.0000E+00~~2.0000E+00~~1.0000E+00~~6.0000E+01
\\						 \tt%\hspace*{1ex}
5.0000E-01~~3.0000E-01~~2.0000E-01~~1.0000E+00~~~~CF1
\\						 \tt\hspace*{-2.5ex}
-1.9570E-01~~6.5000E-01~~4.0000E-01~~8.5500E-02~~1.0000E+00
\\						 \tt%\hspace*{1ex}
1.0000E+00~~2.0000E+00~~1.0000E+00~~2.0000E+00~~0.0000E+00
\end{tabular}\hphantom{AAAAAA}\end{equation}
The first line in every group is also ignored by the code.
It gives the ordinal number of the group (or of the corresponding pencil
beam) and is useful for manual creating or editing the configuration
file only. 
The other parameters have the following significance

\noindent
{\bf Table 4.31. Parameters describing each pencil beam}\\[2ex]
\begin{tabular}{| c | c | l |}
\hline
\bf No.&\bf\hspace*{-1pt}Units\hspace*{-1pt}&\bf Description\\
\hline
1  &[MW]& Pencil beam power					\\
2  &--& Counter injection fraction (0 for co-, 1 for counter-injection)\\
3  &--& Mass of the beam ions in the proton mass (can be 1,2,3)\\
4  &--& Charge of the beam ions in the proton charge units	\\
5  &--& Maximum beam energy	\\
6  &--& Power fraction of the full (maximum) energy component 	\\
7  &--& Power fraction of the one half energy component	\\
8  &--& Power fraction of the one third energy component	\\
9  &--& Type of averaging over the fast ion orbits	\\
10 &--& Number of thin beams in the horizontal plane	\\
11 &[m]& Vertical shift in the beam footprint centre	\\
12 &[m]& Maximum major radius in the beam footprint	\\
13 &[m]& Minimum major radius in the beam footprint	\\
14 &--& $\tan\al,$ $\al$ being the angle between the beam axis and the midplane\\
15 &--& Footprint aspect ratio	\\ \hline
16 &--& 
These 4 parameters prescribe the beam power distribution across the\\
17 &--& 
pencil beam. 
The parameters are transferred to the user defined functions\\
18 &--& 
\twd{NBFHZ} and \twd{NBFRY} (see the file \twd{AWD/sbr/nbuser.f}) 
under the names\\
19 &--& \twd{CVER1, CVER2, CHOR1, CHOR2}, respectively	\\ \hline
20 &--& Unused parameter				\\
\hline
\end{tabular}\\[15pt]
The first parameter gives the power injected by each beam source.
This quantity is usually time dependent, therefore, it is convenient to
set it as a variable rather than as a number. 
In the example \rf{NBI_conf}, this variable is $\tt ZRD1\equiv ZRD1X$
and its time evolution can be prescribed in the data file
(Section~5.3.3). 
Similarly, the number of partial thin beams\footnote
{The parameter No.~10 in Table~4.31 gives the number of beams in the
horizontal direction, $N_h.$ The total number of constituent thin beams
is evaluated as $N_{tot}=N_hN_v$ while $N_v$ is calculated by the code
as a number of flux surfaces in the vertical aperture of each pencil
beam in the footprint crossection.
} 
which represent the pencil beam is an adjustable parameter
that can be varied interactively (\twd{CF1} in example \rf{NBI_conf}). 

The 9-th parameter in Table~4.31 also requires more explanations.
This parameter defines how the power of a fast ion is deposited in the
plasma: \\
\hspace*{1em}{\tt $=0$} no averaging (deposition at the birth point),\\
\hspace*{1em}{\tt $=1$} averaging with a finite orbit width,\\
\hspace*{1em}{\tt $=2$} averaging with zero orbit width.\\
\vspace*{-1ex}

The last 10 parameters in Table~4.31 define the geometry of every pencil
beam.
The neutral beam is characterized, first of all, by its ``footprint'',
i.e. the cross-section of the beam with the plane including the major
torus axis and orthogonal to the beam axis projection onto the midplane.
In turn, the footprint represented as a rectangle is given by four
parameters: the rectangle center (parameter No.~11), major radii of two
vertical sides (parameters 12 and 13) and by the aspect ratio that is
defined as ratio of the rectangle height to its width.
Secondly, by the inclination of the beam axis with respect to the
midplane (parameter No.~14), and, thirdly,\footnote
{In the current version of the code, the beam divergence is not taken
into account.}
by the power distribution across the beam axis.
The power distribution is described by two Fortran functions,
\twd{NBFHZ} and \twd{NBFRY}, in the vertical and in the horizontal
cross-sections, respectively. 
Each function has two free parameters which are the first four
parameters in the last line of example \rf{NBI_conf} (or parameters 16,
17, 18, 19 in Table~4.31). 
Both functions are included in the file \twd{AWD/sbr/nbuser.f} and can
be arbitrarily modified by the user to prescribe any power distribution
required. 

The NBI configuration file can either be created by the user from
scratch or adjusted by editing one of the files of this type provided
with the code.
Another option is interactive creating/modification during the
Astra run.
The code enters in the interactive regime whenever\\
\hspace*{1em}-- the NBI configuration file does not exist,\\
\hspace*{1em}-- the parameter {\tt CNB1} takes a negative value,\\
\hspace*{1em}-- the parameter {\tt CNB1} changes its value during the run.\\
In all cases, a dialog window pops up and the user can adjust all input
parameters for all beam sources.
As before, assume that the input data file is \twd{aug10000}, (the full
name of the file being \twd{AWD/exp/aug10000}).
If neither \twd{aug10000.nbi} nor \twd{aug.nbi} exist in the directory
\twd{AWD/exp/} then a new file \twd{AWD/exp/aug10000.nbi} will be
created on exit from the dialog window.
The new (or modified) data file will then be used in the subsequent
simulations. 
The total number of sources is defined as the absolute value
$\lb|\twd{CNB1}\rb|.$ 
Note also that if the actual number of records in the file is larger
than expected due to $\lb|\twd{CNB1}\rb|$ then the excessive records
are ignored.
The opposite case and $\tw{CNB1}=0$ are treated as an error.


The NBI routine provides the following output.

\noindent
{\bf Table 4.32. Output from \twd{NBI} routine}\\[2ex]
\begin{tabular}{| l l l |}					\hline
\bf Parameter	&\bf Units	&\bf Description\\		\hline
{\tt CUBM}	&[MA$\,\rm m^{-2}$]	& NB driven current\\
{\tt CUFI}	&[MA$\,\rm m^{-2}$]	& Current of the fast ions\\
{\tt NIBM}	&[10$^{19}{\rm m}^{-3}$]
		&Density of the suprathermal ions\\
{\tt NNBM1}	&[10$^{19}{\rm m}^{-3}$]	
		&Density of neutrals with the full energy {\tt EBEAM}\\ 
{\tt NNBM2}	&[10$^{19}{\rm m}^{-3}$]
		&Density of neutrals with the half energy {\tt EBEAM}/2\\
{\tt NNBM3}	&[10$^{19}{\rm m}^{-3}$]
		& Density of neutrals with the energy {\tt EBEAM}/3\\
{\tt PBEAM}	&[MW${\rm m}^{-3}$]& Total beam power absorbed\\
{\tt PBLON}	&[$10^{19}$keV${\rm m}^{-3}$]	
		& Longitudinal pressure due to fast ions\\
{\tt PBPER}	&[$10^{19}$keV${\rm m}^{-3}$]	
		& Perpendicular pressure due to fast ions\\
{\tt PEBM}	&[MW${\rm m}^{-3}$]& Beam power absorbed by electrons\\
{\tt PIBM}	&[MW${\rm m}^{-3}$]& Beam power absorbed by ions\\
{\tt SCUBM}	&[kg~s$^{-2}$m$^{-2}$]	
		& Source of the toroidal momentum\\
{\tt SNEBM}	&[$10^{19}$m$^{-3}$s$^{-1}$]	
		& Source of electrons due to NBI\\
{\tt SNIBM1} &$[10^{19}$m$^{-3}$s$^{-1}]$ 
		& Birth rate for ions with the energy \twd{EBEAM}\\
{\tt SNIBM2} &$[10^{19}$m$^{-3}$s$^{-1}]$ 
		& Birth rate for ions with the energy \twd{EBEAM}/2\\
{\tt SNIBM3} &$[10^{19}$m$^{-3}$s$^{-1}]$ 
		& Birth rate for ions with the energy \twd{EBEAM}/3\\
{\tt SNNBM}	&[$10^{19}$m$^{-3}$s$^{-1}$]
		& Source of thermal neutrals\\			\hline
\end{tabular}\\[15pt]
More detailed description of the algorithm and a way of calculating all
quantities in Table~4.32 are described in \cite{Pole}. 

\vfill\eject
